\chapter{帮助与排错}

% === 源文件: help/debugging.md ===
\section{调试}


本页介绍用于流式输出的调试辅助工具，特别是当提供商将推理混入正常文本时。

\subsection{运行时调试覆盖}


在聊天中使用 \texttt{/debug} 设置\textbf{仅运行时}配置覆盖（内存中，不写入磁盘）。
\texttt{/debug} 默认禁用；通过 \texttt{commands.debug: true} 启用。
当你需要切换不常用的设置而不编辑 \texttt{openclaw.json} 时，这非常方便。

示例：

\begin{lstlisting}
/debug show
/debug set messages.responsePrefix="[openclaw]"
/debug unset messages.responsePrefix
/debug reset
\end{lstlisting}


\texttt{/debug reset} 清除所有覆盖并返回到磁盘上的配置。

\subsection{Gateway 网关监视模式}


为了快速迭代，在文件监视器下运行 Gateway 网关：

\begin{lstlisting}[language=bash]
pnpm gateway:watch --force
\end{lstlisting}


这映射到：

\begin{lstlisting}[language=bash]
tsx watch src/entry.ts gateway --force
\end{lstlisting}


在 \texttt{gateway:watch} 后添加任何 Gateway 网关 CLI 标志，它们将在每次重启时传递。

\subsection{Dev 配置文件 + dev Gateway 网关（--dev）}


使用 dev 配置文件来隔离状态，并启动一个安全、可丢弃的调试设置。有\textbf{两个} \texttt{--dev} 标志：

\begin{itemize}
  \item \textbf{全局 \texttt{--dev}（配置文件）：} 将状态隔离到 \texttt{\textasciitilde{}/.openclaw-dev} 下，并将 Gateway 网关端口默认为 \texttt{19001}（派生端口随之移动）。
  \item \textbf{\texttt{gateway --dev}：告诉 Gateway 网关在缺失时自动创建默认配置 + 工作区}（并跳过 BOOTSTRAP.md）。
\end{itemize}


推荐流程（dev 配置文件 + dev 引导）：

\begin{lstlisting}[language=bash]
pnpm gateway:dev
OPENCLAW_PROFILE=dev openclaw tui
\end{lstlisting}


如果你还没有全局安装，请通过 \texttt{pnpm openclaw ...} 运行 CLI。

这会执行：

\begin{enumerate}
  \item \textbf{配置文件隔离}（全局 \texttt{--dev}）
\end{enumerate}

\begin{itemize}
  \item \texttt{OPENCLAW\_PROFILE=dev}
  \item \texttt{OPENCLAW\_STATE\_DIR=\textasciitilde{}/.openclaw-dev}
  \item \texttt{OPENCLAW\_CONFIG\_PATH=\textasciitilde{}/.openclaw-dev/openclaw.json}
  \item \texttt{OPENCLAW\_GATEWAY\_PORT=19001}（浏览器/画布相应移动）
\end{itemize}


\begin{enumerate}
  \item \textbf{Dev 引导}（\texttt{gateway --dev}）
\end{enumerate}

\begin{itemize}
  \item 如果缺失则写入最小配置（\texttt{gateway.mode=local}，绑定 loopback）。
  \item 将 \texttt{agent.workspace} 设置为 dev 工作区。
  \item 设置 \texttt{agent.skipBootstrap=true}（无 BOOTSTRAP.md）。
  \item 如果缺失则填充工作区文件：
\end{itemize}

\texttt{AGENTS.md}、\texttt{SOUL.md}、\texttt{TOOLS.md}、\texttt{IDENTITY.md}、\texttt{USER.md}、\texttt{HEARTBEAT.md}。
\begin{itemize}
  \item 默认身份：\textbf{C3‑PO}（礼仪机器人）。
  \item 在 dev 模式下跳过渠道提供商（\texttt{OPENCLAW\_SKIP\_CHANNELS=1}）。
\end{itemize}


重置流程（全新开始）：

\begin{lstlisting}[language=bash]
pnpm gateway:dev:reset
\end{lstlisting}


注意：\texttt{--dev} 是\textbf{全局}配置文件标志，会被某些运行器吞掉。
如果你需要明确拼写，请使用环境变量形式：

\begin{lstlisting}[language=bash]
OPENCLAW_PROFILE=dev openclaw gateway --dev --reset
\end{lstlisting}


\texttt{--reset} 清除配置、凭证、会话和 dev 工作区（使用 \texttt{trash}，而非 \texttt{rm}），然后重新创建默认的 dev 设置。

提示：如果非 dev Gateway 网关已在运行（launchd/systemd），请先停止它：

\begin{lstlisting}[language=bash]
openclaw gateway stop
\end{lstlisting}


\subsection{原始流日志（OpenClaw）}


OpenClaw 可以在任何过滤/格式化之前记录\textbf{原始助手流}。
这是查看推理是否作为纯文本增量到达（或作为单独的思考块）的最佳方式。

通过 CLI 启用：

\begin{lstlisting}[language=bash]
pnpm gateway:watch --force --raw-stream
\end{lstlisting}


可选路径覆盖：

\begin{lstlisting}[language=bash]
pnpm gateway:watch --force --raw-stream --raw-stream-path ~/.openclaw/logs/raw-stream.jsonl
\end{lstlisting}


等效环境变量：

\begin{lstlisting}[language=bash]
OPENCLAW_RAW_STREAM=1
OPENCLAW_RAW_STREAM_PATH=~/.openclaw/logs/raw-stream.jsonl
\end{lstlisting}


默认文件：

\texttt{\textasciitilde{}/.openclaw/logs/raw-stream.jsonl}

\subsection{原始块日志（pi-mono）}


要在解析为块之前捕获\textbf{原始 OpenAI 兼容块}，pi-mono 暴露了一个单独的日志记录器：

\begin{lstlisting}[language=bash]
PI_RAW_STREAM=1
\end{lstlisting}


可选路径：

\begin{lstlisting}[language=bash]
PI_RAW_STREAM_PATH=~/.pi-mono/logs/raw-openai-completions.jsonl
\end{lstlisting}


默认文件：

\texttt{\textasciitilde{}/.pi-mono/logs/raw-openai-completions.jsonl}

\begin{quote}
注意：这仅由使用 pi-mono 的 \texttt{openai-completions} 提供商的进程发出。
\end{quote}


\subsection{安全注意事项}


\begin{itemize}
  \item 原始流日志可能包含完整提示、工具输出和用户数据。
  \item 保持日志在本地并在调试后删除它们。
  \item 如果你分享日志，请先清除密钥和个人身份信息。
\end{itemize}



\clearpage

% === 源文件: help/environment.md ===
\section{环境变量}


OpenClaw 从多个来源拉取环境变量。规则是\textbf{永不覆盖现有值}。

\subsection{优先级（从高到低）}


\begin{enumerate}
  \item \textbf{进程环境}（Gateway 网关进程从父 shell/守护进程已有的内容）。
  \item \textbf{当前工作目录中的 \texttt{.env}}（dotenv 默认；不覆盖）。
  \item \textbf{全局 \texttt{.env}} 位于 \texttt{\textasciitilde{}/.openclaw/.env}（即 \texttt{\$OPENCLAW\_STATE\_DIR/.env}；不覆盖）。
  \item \textbf{配置 \texttt{env} 块} 位于 \texttt{\textasciitilde{}/.openclaw/openclaw.json}（仅在缺失时应用）。
  \item \textbf{可选的登录 shell 导入}（\texttt{env.shellEnv.enabled} 或 \texttt{OPENCLAW\_LOAD\_SHELL\_ENV=1}），仅对缺失的预期键名应用。
\end{enumerate}


如果配置文件完全缺失，步骤 4 将被跳过；如果启用了 shell 导入，它仍会运行。

\subsection{配置 env 块}


两种等效方式设置内联环境变量（都是非覆盖的）：

\begin{lstlisting}[language=json]
{
  env: {
    OPENROUTER_API_KEY: "sk-or-...",
    vars: {
      GROQ_API_KEY: "gsk-...",
    },
  },
}
\end{lstlisting}


\subsection{Shell 环境导入}


\texttt{env.shellEnv} 运行你的登录 shell 并仅导入\textbf{缺失的}预期键名：

\begin{lstlisting}[language=json]
{
  env: {
    shellEnv: {
      enabled: true,
      timeoutMs: 15000,
    },
  },
}
\end{lstlisting}


环境变量等效项：

\begin{itemize}
  \item \texttt{OPENCLAW\_LOAD\_SHELL\_ENV=1}
  \item \texttt{OPENCLAW\_SHELL\_ENV\_TIMEOUT\_MS=15000}
\end{itemize}


\subsection{配置中的环境变量替换}


你可以使用 \texttt{\$\{VAR\_NAME\}} 语法在配置字符串值中直接引用环境变量：

\begin{lstlisting}[language=json]
{
  models: {
    providers: {
      "vercel-gateway": {
        apiKey: "${VERCEL_GATEWAY_API_KEY}",
      },
    },
  },
}
\end{lstlisting}


完整详情参见\href{/gateway/configuration\#env-var-substitution-in-config}{配置：环境变量替换}。

\subsection{相关内容}


\begin{itemize}
  \item \href{/gateway/configuration}{Gateway 网关配置}
  \item \href{/help/faq\#env-vars-and-env-loading}{常见问题：环境变量和 .env 加载}
  \item \href{/concepts/models}{模型概述}
\end{itemize}



\clearpage

% === 源文件: help/faq.md ===
\section{常见问题}


快速解答及针对实际部署场景（本地开发、VPS、多智能体、OAuth/API 密钥、模型故障转移）的深入故障排除。运行时诊断请参阅\href{/gateway/troubleshooting}{故障排除}。完整配置参考请参阅\href{/gateway/configuration}{配置}。

\subsection{目录}


\begin{itemize}
  \item \href{\#quick-start-and-firstrun-setup}{快速开始与首次运行设置}
  \item \href{\#im-stuck-whats-the-fastest-way-to-get-unstuck}{我卡住了，最快的排障方法是什么？}
  \item \href{\#whats-the-recommended-way-to-install-and-set-up-openclaw}{安装和设置 OpenClaw 的推荐方式是什么？}
  \item \href{\#how-do-i-open-the-dashboard-after-onboarding}{新手引导后如何打开仪表板？}
  \item \href{\#how-do-i-authenticate-the-dashboard-token-on-localhost-vs-remote}{如何在本地和远程环境中验证仪表板（令牌）？}
  \item \href{\#what-runtime-do-i-need}{我需要什么运行时？}
  \item \href{\#does-it-run-on-raspberry-pi}{能在 Raspberry Pi 上运行吗？}
  \item \href{\#any-tips-for-raspberry-pi-installs}{Raspberry Pi 安装有什么建议？}
  \item \href{\#it-is-stuck-on-wake-up-my-friend-onboarding-will-not-hatch-what-now}{卡在 "wake up my friend" / 新手引导无法启动，怎么办？}
  \item \href{\#can-i-migrate-my-setup-to-a-new-machine-mac-mini-without-redoing-onboarding}{能否将我的设置迁移到新机器（Mac mini）而不重新进行新手引导？}
  \item \href{\#where-do-i-see-whats-new-in-the-latest-version}{在哪里查看最新版本的更新内容？}
  \item \href{\#i-cant-access-docsopenclawai-ssl-error-what-now}{无法访问 docs.openclaw.ai（SSL 错误），怎么办？}
  \item \href{\#whats-the-difference-between-stable-and-beta}{stable 和 beta 有什么区别？}
  \item \href{\#how-do-i-install-the-beta-version-and-whats-the-difference-between-beta-and-dev}{如何安装 beta 版本，beta 和 dev 有什么区别？}
  \item \href{\#how-do-i-try-the-latest-bits}{如何试用最新代码？}
  \item \href{\#how-long-does-install-and-onboarding-usually-take}{安装和新手引导通常需要多长时间？}
  \item \href{\#installer-stuck-how-do-i-get-more-feedback}{安装程序卡住了？如何获取更多反馈？}
  \item \href{\#windows-install-says-git-not-found-or-openclaw-not-recognized}{Windows 安装提示找不到 git 或无法识别 openclaw}
  \item \href{\#the-docs-didnt-answer-my-question-how-do-i-get-a-better-answer}{文档没有解答我的问题——如何获得更好的答案？}
  \item \href{\#how-do-i-install-openclaw-on-linux}{如何在 Linux 上安装 OpenClaw？}
  \item \href{\#how-do-i-install-openclaw-on-a-vps}{如何在 VPS 上安装 OpenClaw？}
  \item \href{\#where-are-the-cloudvps-install-guides}{云/VPS 安装指南在哪里？}
  \item \href{\#can-i-ask-openclaw-to-update-itself}{可以让 OpenClaw 自行更新吗？}
  \item \href{\#what-does-the-onboarding-wizard-actually-do}{新手引导向导具体做了什么？}
  \item \href{\#do-i-need-a-claude-or-openai-subscription-to-run-this}{运行 OpenClaw 需要 Claude 或 OpenAI 订阅吗？}
  \item \href{\#can-i-use-claude-max-subscription-without-an-api-key}{能否使用 Claude Max 订阅而不需要 API 密钥？}
  \item \href{\#how-does-anthropic-setuptoken-auth-work}{Anthropic "setup-token" 认证如何工作？}
  \item \href{\#where-do-i-find-an-anthropic-setuptoken}{在哪里获取 Anthropic setup-token？}
  \item \href{\#do-you-support-claude-subscription-auth-claude-code-oauth}{是否支持 Claude 订阅认证（Claude Code OAuth）？}
  \item \href{\#why-am-i-seeing-http-429-ratelimiterror-from-anthropic}{为什么我看到 \texttt{HTTP 429: rate\_limit\_error}（来自 Anthropic）？}
  \item \href{\#is-aws-bedrock-supported}{支持 AWS Bedrock 吗？}
  \item \href{\#how-does-codex-auth-work}{Codex 认证如何工作？}
  \item \href{\#do-you-support-openai-subscription-auth-codex-oauth}{是否支持 OpenAI 订阅认证（Codex OAuth）？}
  \item \href{\#how-do-i-set-up-gemini-cli-oauth}{如何设置 Gemini CLI OAuth？}
  \item \href{\#is-a-local-model-ok-for-casual-chats}{本地模型适合日常聊天吗？}
  \item \href{\#how-do-i-keep-hosted-model-traffic-in-a-specific-region}{如何将托管模型流量限制在特定区域？}
  \item \href{\#do-i-have-to-buy-a-mac-mini-to-install-this}{我必须购买 Mac Mini 才能安装吗？}
  \item \href{\#do-i-need-a-mac-mini-for-imessage-support}{iMessage 支持需要 Mac mini 吗？}
  \item \href{\#if-i-buy-a-mac-mini-to-run-openclaw-can-i-connect-it-to-my-macbook-pro}{如果我买了 Mac mini 运行 OpenClaw，能连接到我的 MacBook Pro 吗？}
  \item \href{\#can-i-use-bun}{可以使用 Bun 吗？}
  \item \href{\#telegram-what-goes-in-allowfrom}{Telegram：\texttt{allowFrom} 填什么？}
  \item \href{\#can-multiple-people-use-one-whatsapp-number-with-different-openclaw-instances}{多人能否使用同一个 WhatsApp 号码配合不同的 OpenClaw 实例？}
  \item \href{\#can-i-run-a-fast-chat-agent-and-an-opus-for-coding-agent}{能否同时运行一个“快速聊天”智能体和一个“用 Opus 编程”的智能体？}
  \item \href{\#does-homebrew-work-on-linux}{Homebrew 在 Linux 上可用吗？}
  \item \href{\#whats-the-difference-between-the-hackable-git-install-and-npm-install}{可编辑（git）安装和 npm 安装有什么区别？}
  \item \href{\#can-i-switch-between-npm-and-git-installs-later}{之后可以在 npm 和 git 安装之间切换吗？}
  \item \href{\#should-i-run-the-gateway-on-my-laptop-or-a-vps}{应该在笔记本电脑还是 VPS 上运行 Gateway 网关？}
  \item \href{\#how-important-is-it-to-run-openclaw-on-a-dedicated-machine}{在专用机器上运行 OpenClaw 有多重要？}
  \item \href{\#what-are-the-minimum-vps-requirements-and-recommended-os}{VPS 的最低要求和推荐操作系统是什么？}
  \item \href{\#can-i-run-openclaw-in-a-vm-and-what-are-the-requirements}{可以在虚拟机中运行 OpenClaw 吗？有什么要求？}
  \item \href{\#what-is-openclaw}{什么是 OpenClaw？}
  \item \href{\#what-is-openclaw-in-one-paragraph}{用一段话描述 OpenClaw？}
  \item \href{\#whats-the-value-proposition}{价值主张是什么？}
  \item \href{\#i-just-set-it-up-what-should-i-do-first}{刚设置好，应该先做什么？}
  \item \href{\#what-are-the-top-five-everyday-use-cases-for-openclaw}{OpenClaw 日常最常用的五个场景是什么？}
  \item \href{\#can-openclaw-help-with-lead-gen-outreach-ads-and-blogs-for-a-saas}{OpenClaw 能否帮助 SaaS 进行获客、外联、广告和博客？}
  \item \href{\#what-are-the-advantages-vs-claude-code-for-web-development}{相比 Claude Code，在 Web 开发方面有什么优势？}
  \item \href{\#skills-and-automation}{Skills 与自动化}
  \item \href{\#how-do-i-customize-skills-without-keeping-the-repo-dirty}{如何自定义 Skills 而不弄脏仓库？}
  \item \href{\#can-i-load-skills-from-a-custom-folder}{可以从自定义文件夹加载 Skills 吗？}
  \item \href{\#how-can-i-use-different-models-for-different-tasks}{如何为不同任务使用不同模型？}
  \item \href{\#the-bot-freezes-while-doing-heavy-work-how-do-i-offload-that}{机器人在执行繁重工作时卡住了，如何卸载任务？}
  \item \href{\#cron-or-reminders-do-not-fire-what-should-i-check}{定时任务或提醒没有触发，应该检查什么？}
  \item \href{\#how-do-i-install-skills-on-linux}{如何在 Linux 上安装 Skills？}
  \item \href{\#can-openclaw-run-tasks-on-a-schedule-or-continuously-in-the-background}{OpenClaw 能否按计划或在后台持续运行任务？}
  \item \href{\#can-i-run-applemacosonly-skills-from-linux}{能否从 Linux 运行仅限 Apple/macOS 的 Skills？}
  \item \href{\#do-you-have-a-notion-or-heygen-integration}{有 Notion 或 HeyGen 集成吗？}
  \item \href{\#how-do-i-install-the-chrome-extension-for-browser-takeover}{如何安装用于浏览器接管的 Chrome 扩展？}
  \item \href{\#sandboxing-and-memory}{沙箱与记忆}
  \item \href{\#is-there-a-dedicated-sandboxing-doc}{有专门的沙箱文档吗？}
  \item \href{\#how-do-i-bind-a-host-folder-into-the-sandbox}{如何将主机文件夹绑定到沙箱中？}
  \item \href{\#how-does-memory-work}{记忆是如何工作的？}
  \item \href{\#memory-keeps-forgetting-things-how-do-i-make-it-stick}{记忆总是遗忘，如何让它持久保存？}
  \item \href{\#does-memory-persist-forever-what-are-the-limits}{记忆是否永久保留？有什么限制？}
  \item \href{\#does-semantic-memory-search-require-an-openai-api-key}{语义记忆搜索需要 OpenAI API 密钥吗？}
  \item \href{\#where-things-live-on-disk}{磁盘上的文件位置}
  \item \href{\#is-all-data-used-with-openclaw-saved-locally}{OpenClaw 使用的所有数据都保存在本地吗？}
  \item \href{\#where-does-openclaw-store-its-data}{OpenClaw 将数据存储在哪里？}
  \item \href{\#where-should-agentsmd-soulmd-usermd-memorymd-live}{AGENTS.md / SOUL.md / USER.md / MEMORY.md 应该放在哪里？}
  \item \href{\#whats-the-recommended-backup-strategy}{推荐的备份策略是什么？}
  \item \href{\#how-do-i-completely-uninstall-openclaw}{如何完全卸载 OpenClaw？}
  \item \href{\#can-agents-work-outside-the-workspace}{智能体可以在工作区外工作吗？}
  \item \href{\#im-in-remote-mode-where-is-the-session-store}{我处于远程模式——会话存储在哪里？}
  \item \href{\#config-basics}{配置基础}
  \item \href{\#what-format-is-the-config-where-is-it}{配置文件是什么格式？在哪里？}
  \item \href{\#i-set-gatewaybind-lan-or-tailnet-and-now-nothing-listens-the-ui-says-unauthorized}{我设置了 \texttt{gateway.bind: "lan"}（或 \texttt{"tailnet"}），现在什么都监听不了 / UI 显示未授权}
  \item \href{\#why-do-i-need-a-token-on-localhost-now}{为什么现在在 localhost 也需要令牌？}
  \item \href{\#do-i-have-to-restart-after-changing-config}{更改配置后需要重启吗？}
  \item \href{\#how-do-i-enable-web-search-and-web-fetch}{如何启用网络搜索（和网页抓取）？}
  \item \href{\#configapply-wiped-my-config-how-do-i-recover-and-avoid-this}{config.apply 清空了我的配置，如何恢复和避免？}
  \item \href{\#how-do-i-run-a-central-gateway-with-specialized-workers-across-devices}{如何运行一个中心 Gateway 网关配合跨设备的专用工作节点？}
  \item \href{\#can-the-openclaw-browser-run-headless}{OpenClaw 浏览器可以无头运行吗？}
  \item \href{\#how-do-i-use-brave-for-browser-control}{如何使用 Brave 进行浏览器控制？}
  \item \href{\#remote-gateways-nodes}{远程 Gateway 网关与节点}
  \item \href{\#how-do-commands-propagate-between-telegram-the-gateway-and-nodes}{命令如何在 Telegram、Gateway 网关和节点之间传播？}
  \item \href{\#how-can-my-agent-access-my-computer-if-the-gateway-is-hosted-remotely}{如果 Gateway 网关托管在远程，我的智能体如何访问我的电脑？}
  \item \href{\#tailscale-is-connected-but-i-get-no-replies-what-now}{Tailscale 已连接但收不到回复，怎么办？}
  \item \href{\#can-two-openclaw-instances-talk-to-each-other-local-vps}{两个 OpenClaw 实例（本地 + VPS）可以互相通信吗？}
  \item \href{\#do-i-need-separate-vpses-for-multiple-agents}{多个智能体需要独立的 VPS 吗？}
  \item \href{\#is-there-a-benefit-to-using-a-node-on-my-personal-laptop-instead-of-ssh-from-a-vps}{在个人笔记本电脑上使用节点而不是从 VPS SSH 有什么好处？}
  \item \href{\#do-nodes-run-a-gateway-service}{节点会运行 Gateway 网关服务吗？}
  \item \href{\#is-there-an-api-rpc-way-to-apply-config}{有 API / RPC 方式来应用配置吗？}
  \item \href{\#whats-a-minimal-sane-config-for-a-first-install}{首次安装的最小“合理”配置是什么？}
  \item \href{\#how-do-i-set-up-tailscale-on-a-vps-and-connect-from-my-mac}{如何在 VPS 上设置 Tailscale 并从 Mac 连接？}
  \item \href{\#how-do-i-connect-a-mac-node-to-a-remote-gateway-tailscale-serve}{如何将 Mac 节点连接到远程 Gateway 网关（Tailscale Serve）？}
  \item \href{\#should-i-install-on-a-second-laptop-or-just-add-a-node}{应该在第二台笔记本上安装还是只添加一个节点？}
  \item \href{\#env-vars-and-env-loading}{环境变量和 .env 加载}
  \item \href{\#how-does-openclaw-load-environment-variables}{OpenClaw 如何加载环境变量？}
  \item \href{\#i-started-the-gateway-via-the-service-and-my-env-vars-disappeared-what-now}{“我通过服务启动了 Gateway 网关，但环境变量消失了。”怎么办？}
  \item \href{\#i-set-copilotgithubtoken-but-models-status-shows-shell-env-off-why}{我设置了 \texttt{COPILOT\_GITHUB\_TOKEN}，但 models status 显示"Shell env: off"，为什么？}
  \item \href{\#sessions-multiple-chats}{会话与多聊天}
  \item \href{\#how-do-i-start-a-fresh-conversation}{如何开始一个新对话？}
  \item \href{\#do-sessions-reset-automatically-if-i-never-send-new}{如果我从不发送 \texttt{/new}，会话会自动重置吗？}
  \item \href{\#is-there-a-way-to-make-a-team-of-openclaw-instances-one-ceo-and-many-agents}{能否创建一个 OpenClaw 实例团队——一个 CEO 和多个智能体？}
  \item \href{\#why-did-context-get-truncated-midtask-how-do-i-prevent-it}{为什么上下文在任务中途被截断了？如何防止？}
  \item \href{\#how-do-i-completely-reset-openclaw-but-keep-it-installed}{如何完全重置 OpenClaw 但保留安装？}
  \item \href{\#im-getting-context-too-large-errors-how-do-i-reset-or-compact}{我遇到了"context too large"错误——如何重置或压缩？}
  \item \href{\#why-am-i-seeing-llm-request-rejected-messagesncontentxtooluseinput-field-required}{为什么我看到"LLM request rejected: messages.N.content.X.tool\_use.input: Field required"？}
  \item \href{\#why-am-i-getting-heartbeat-messages-every-30-minutes}{为什么每 30 分钟收到一次心跳消息？}
  \item \href{\#do-i-need-to-add-a-bot-account-to-a-whatsapp-group}{需要在 WhatsApp 群组中添加“机器人账号”吗？}
  \item \href{\#how-do-i-get-the-jid-of-a-whatsapp-group}{如何获取 WhatsApp 群组的 JID？}
  \item \href{\#why-doesnt-openclaw-reply-in-a-group}{为什么 OpenClaw 不在群组中回复？}
  \item \href{\#do-groupsthreads-share-context-with-dms}{群组/线程与私聊共享上下文吗？}
  \item \href{\#how-many-workspaces-and-agents-can-i-create}{可以创建多少个工作区和智能体？}
  \item \href{\#can-i-run-multiple-bots-or-chats-at-the-same-time-slack-and-how-should-i-set-that-up}{可以同时运行多个机器人或聊天（Slack）吗？应该如何设置？}
  \item \href{\#models-defaults-selection-aliases-switching}{模型：默认值、选择、别名、切换}
  \item \href{\#what-is-the-default-model}{什么是“默认模型”？}
  \item \href{\#what-model-do-you-recommend}{推荐什么模型？}
  \item \href{\#how-do-i-switch-models-without-wiping-my-config}{如何在不清空配置的情况下切换模型？}
  \item \href{\#can-i-use-selfhosted-models-llamacpp-vllm-ollama}{可以使用自托管模型（llama.cpp、vLLM、Ollama）吗？}
  \item \href{\#what-do-openclaw-flawd-and-krill-use-for-models}{OpenClaw、Flawd 和 Krill 使用什么模型？}
  \item \href{\#how-do-i-switch-models-on-the-fly-without-restarting}{如何在运行中切换模型（无需重启）？}
  \item \href{\#can-i-use-gpt-52-for-daily-tasks-and-codex-52-for-coding}{能否日常任务用 GPT 5.2，编程用 Codex 5.2？}
  \item \href{\#why-do-i-see-model-is-not-allowed-and-then-no-reply}{为什么我看到"Model … is not allowed"然后没有回复？}
  \item \href{\#why-do-i-see-unknown-model-minimaxminimaxm21}{为什么我看到"Unknown model: minimax/MiniMax-M2.1"？}
  \item \href{\#can-i-use-minimax-as-my-default-and-openai-for-complex-tasks}{能否将 MiniMax 设为默认，复杂任务用 OpenAI？}
  \item \href{\#are-opus-sonnet-gpt-builtin-shortcuts}{opus / sonnet / gpt 是内置快捷方式吗？}
  \item \href{\#how-do-i-defineoverride-model-shortcuts-aliases}{如何定义/覆盖模型快捷方式（别名）？}
  \item \href{\#how-do-i-add-models-from-other-providers-like-openrouter-or-zai}{如何添加其他提供商（如 OpenRouter 或 Z.AI）的模型？}
  \item \href{\#model-failover-and-all-models-failed}{模型故障转移与"All models failed"}
  \item \href{\#how-does-failover-work}{故障转移是如何工作的？}
  \item \href{\#what-does-this-error-mean}{这个错误是什么意思？}
  \item \href{\#fix-checklist-for-no-credentials-found-for-profile-anthropicdefault}{\texttt{No credentials found for profile "anthropic:default"} 的修复清单}
  \item \href{\#why-did-it-also-try-google-gemini-and-fail}{为什么还尝试了 Google Gemini 并且失败了？}
  \item \href{\#auth-profiles-what-they-are-and-how-to-manage-them}{认证配置文件：概念和管理方式}
  \item \href{\#what-is-an-auth-profile}{什么是认证配置文件？}
  \item \href{\#what-are-typical-profile-ids}{典型的配置文件 ID 有哪些？}
  \item \href{\#can-i-control-which-auth-profile-is-tried-first}{可以控制首先尝试哪个认证配置文件吗？}
  \item \href{\#oauth-vs-api-key-whats-the-difference}{OAuth 与 API 密钥：有什么区别？}
  \item \href{\#gateway-ports-already-running-and-remote-mode}{Gateway 网关：端口、“已在运行”和远程模式}
  \item \href{\#what-port-does-the-gateway-use}{Gateway 网关使用什么端口？}
  \item \href{\#why-does-openclaw-gateway-status-say-runtime-running-but-rpc-probe-failed}{为什么 \texttt{openclaw gateway status} 显示 \texttt{Runtime: running} 但 \texttt{RPC probe: failed}？}
  \item \href{\#why-does-openclaw-gateway-status-show-config-cli-and-config-service-different}{为什么 \texttt{openclaw gateway status} 显示 \texttt{Config (cli)} 和 \texttt{Config (service)} 不同？}
  \item \href{\#what-does-another-gateway-instance-is-already-listening-mean}{"another gateway instance is already listening"是什么意思？}
  \item \href{\#how-do-i-run-openclaw-in-remote-mode-client-connects-to-a-gateway-elsewhere}{如何以远程模式运行 OpenClaw（客户端连接到其他位置的 Gateway 网关）？}
  \item \href{\#the-control-ui-says-unauthorized-or-keeps-reconnecting-what-now}{控制 UI 显示"unauthorized"（或持续重连），怎么办？}
  \item \href{\#i-set-gatewaybind-tailnet-but-it-cant-bind-nothing-listens}{我设置了 \texttt{gateway.bind: "tailnet"} 但无法绑定 / 什么都没监听}
  \item \href{\#can-i-run-multiple-gateways-on-the-same-host}{可以在同一主机上运行多个 Gateway 网关吗？}
  \item \href{\#what-does-invalid-handshake-code-1008-mean}{"invalid handshake" / code 1008 是什么意思？}
  \item \href{\#logging-and-debugging}{日志与调试}
  \item \href{\#where-are-logs}{日志在哪里？}
  \item \href{\#how-do-i-startstoprestart-the-gateway-service}{如何启动/停止/重启 Gateway 网关服务？}
  \item \href{\#i-closed-my-terminal-on-windows-how-do-i-restart-openclaw}{我在 Windows 上关闭了终端——如何重启 OpenClaw？}
  \item \href{\#the-gateway-is-up-but-replies-never-arrive-what-should-i-check}{Gateway 网关已启动但回复始终不到达，应该检查什么？}
  \item \href{\#disconnected-from-gateway-no-reason-what-now}{"Disconnected from gateway: no reason"——怎么办？}
  \item \href{\#telegram-setmycommands-fails-with-network-errors-what-should-i-check}{Telegram setMyCommands 因网络错误失败，应该检查什么？}
  \item \href{\#tui-shows-no-output-what-should-i-check}{TUI 没有输出，应该检查什么？}
  \item \href{\#how-do-i-completely-stop-then-start-the-gateway}{如何完全停止然后启动 Gateway 网关？}
  \item \href{\#eli5-openclaw-gateway-restart-vs-openclaw-gateway}{通俗解释：\texttt{openclaw gateway restart} 与 \texttt{openclaw gateway}}
  \item \href{\#whats-the-fastest-way-to-get-more-details-when-something-fails}{出现故障时获取更多详情的最快方法是什么？}
  \item \href{\#media-attachments}{媒体与附件}
  \item \href{\#my-skill-generated-an-imagepdf-but-nothing-was-sent}{我的 Skills 生成了图片/PDF，但什么都没发送}
  \item \href{\#security-and-access-control}{安全与访问控制}
  \item \href{\#is-it-safe-to-expose-openclaw-to-inbound-dms}{将 OpenClaw 暴露给入站私信安全吗？}
  \item \href{\#is-prompt-injection-only-a-concern-for-public-bots}{提示注入只对公开机器人有影响吗？}
  \item \href{\#should-my-bot-have-its-own-email-github-account-or-phone-number}{我的机器人应该有自己的邮箱、GitHub 账户或电话号码吗？}
  \item \href{\#can-i-give-it-autonomy-over-my-text-messages-and-is-that-safe}{我能让它自主管理我的短信吗？这安全吗？}
  \item \href{\#can-i-use-cheaper-models-for-personal-assistant-tasks}{个人助理任务可以使用更便宜的模型吗？}
  \item \href{\#i-ran-start-in-telegram-but-didnt-get-a-pairing-code}{我在 Telegram 中运行了 \texttt{/start} 但没收到配对码}
  \item \href{\#whatsapp-will-it-message-my-contacts-how-does-pairing-work}{WhatsApp：会给我的联系人发消息吗？配对如何工作？}
  \item \href{\#chat-commands-aborting-tasks-and-it-wont-stop}{聊天命令、中止任务和“停不下来”}
  \item \href{\#how-do-i-stop-internal-system-messages-from-showing-in-chat}{如何阻止内部系统消息显示在聊天中？}
  \item \href{\#how-do-i-stopcancel-a-running-task}{如何停止/取消正在运行的任务？}
  \item \href{\#how-do-i-send-a-discord-message-from-telegram-crosscontext-messaging-denied}{如何从 Telegram 发送 Discord 消息？（"Cross-context messaging denied"）}
  \item \href{\#why-does-it-feel-like-the-bot-ignores-rapidfire-messages}{为什么感觉机器人“忽略”了快速连发的消息？}
\end{itemize}


\subsection{出问题后的最初六十秒}


\begin{enumerate}
  \item \textbf{快速状态（首先检查）}
\end{enumerate}


\begin{lstlisting}[language=bash]
   openclaw status
\end{lstlisting}


快速本地摘要：操作系统 + 更新、Gateway 网关/服务可达性、智能体/会话、提供商配置 + 运行时问题（Gateway 网关可达时）。

\begin{enumerate}
  \item \textbf{可粘贴的报告（可安全分享）}
\end{enumerate}


\begin{lstlisting}[language=bash]
   openclaw status --all
\end{lstlisting}


只读诊断，附带日志尾部（令牌已脱敏）。

\begin{enumerate}
  \item \textbf{守护进程 + 端口状态}
\end{enumerate}


\begin{lstlisting}[language=bash]
   openclaw gateway status
\end{lstlisting}


显示 supervisor 运行状态与 RPC 可达性、探测目标 URL，以及服务可能使用的配置。

\begin{enumerate}
  \item \textbf{深度探测}
\end{enumerate}


\begin{lstlisting}[language=bash]
   openclaw status --deep
\end{lstlisting}


运行 Gateway 网关健康检查 + 提供商探测（需要可达的 Gateway 网关）。参阅\href{/gateway/health}{健康检查}。

\begin{enumerate}
  \item \textbf{跟踪最新日志}
\end{enumerate}


\begin{lstlisting}[language=bash]
   openclaw logs --follow
\end{lstlisting}


如果 RPC 不可用，回退到：

\begin{lstlisting}[language=bash]
   tail -f "$(ls -t /tmp/openclaw/openclaw-*.log | head -1)"
\end{lstlisting}


文件日志与服务日志是分开的；参阅\href{/logging}{日志}和\href{/gateway/troubleshooting}{故障排除}。

\begin{enumerate}
  \item \textbf{运行 doctor（修复）}
\end{enumerate}


\begin{lstlisting}[language=bash]
   openclaw doctor
\end{lstlisting}


修复/迁移配置/状态 + 运行健康检查。参阅 \href{/gateway/doctor}{Doctor}。

\begin{enumerate}
  \item \textbf{Gateway 网关快照}
\begin{lstlisting}[language=bash]
   openclaw health --json
   openclaw health --verbose   # 出错时显示目标 URL + 配置路径
\end{lstlisting}

\end{enumerate}

向运行中的 Gateway 网关请求完整快照（仅 WS）。参阅\href{/gateway/health}{健康检查}。

\subsection{快速开始与首次运行设置}


\subsubsection{我卡住了，最快的排障方法是什么}


使用能\textbf{看到你机器}的本地 AI 智能体。这比在 Discord 上提问有效得多，因为大多数“卡住了”的情况都是\textbf{本地配置或环境问题}，远程帮助者无法检查。

\begin{itemize}
  \item \textbf{Claude Code}：https://www.anthropic.com/claude-code/
  \item \textbf{OpenAI Codex}：https://openai.com/codex/
\end{itemize}


这些工具可以读取仓库、运行命令、检查日志，并帮助修复你的机器级别设置（PATH、服务、权限、认证文件）。通过可编辑（git）安装提供\textbf{完整源代码}：

\begin{lstlisting}[language=bash]
curl -fsSL https://openclaw.ai/install.sh | bash -s -- --install-method git
\end{lstlisting}


这会从 \textbf{git checkout} 安装 OpenClaw，这样智能体可以读取代码 + 文档，并推理你正在运行的确切版本。你可以随时通过不带 \texttt{--install-method git} 重新运行安装程序切回稳定版。

提示：要求智能体\textbf{计划并监督}修复（逐步进行），然后只执行必要的命令。这样改动较小，更容易审查。

如果你发现了真正的 bug 或修复方案，请提交 GitHub issue 或发送 PR：
https://github.com/openclaw/openclaw/issues
https://github.com/openclaw/openclaw/pulls

从以下命令开始（在寻求帮助时分享输出）：

\begin{lstlisting}[language=bash]
openclaw status
openclaw models status
openclaw doctor
\end{lstlisting}


它们的作用：

\begin{itemize}
  \item \texttt{openclaw status}：Gateway 网关/智能体健康状况 + 基本配置的快速快照。
  \item \texttt{openclaw models status}：检查提供商认证 + 模型可用性。
  \item \texttt{openclaw doctor}：验证并修复常见的配置/状态问题。
\end{itemize}


其他有用的 CLI 检查：\texttt{openclaw status --all}、\texttt{openclaw logs --follow}、
\texttt{openclaw gateway status}、\texttt{openclaw health --verbose}。

快速调试流程：\href{\#first-60-seconds-if-somethings-broken}{出问题后的最初六十秒}。
安装文档：\href{/install}{安装}、\href{/install/installer}{安装程序标志}、\href{/install/updating}{更新}。

\subsubsection{安装和设置 OpenClaw 的推荐方式是什么}


仓库推荐从源码运行并使用新手引导向导：

\begin{lstlisting}[language=bash]
curl -fsSL https://openclaw.ai/install.sh | bash
openclaw onboard --install-daemon
\end{lstlisting}


向导还可以自动构建 UI 资源。新手引导后，通常在端口 \textbf{18789} 上运行 Gateway 网关。

从源码安装（贡献者/开发者）：

\begin{lstlisting}[language=bash]
git clone https://github.com/openclaw/openclaw.git
cd openclaw
pnpm install
pnpm build
pnpm ui:build # 首次运行时自动安装 UI 依赖
openclaw onboard
\end{lstlisting}


如果你还没有全局安装，通过 \texttt{pnpm openclaw onboard} 运行。

\subsubsection{新手引导后如何打开仪表板}


向导现在会在新手引导完成后立即使用带令牌的仪表板 URL 打开浏览器，并在摘要中打印完整链接（带令牌）。保持该标签页打开；如果没有自动启动，请在同一台机器上复制/粘贴打印的 URL。令牌保持在本地主机上——不会从浏览器获取任何内容。

\subsubsection{如何在本地和远程环境中验证仪表板令牌}


\textbf{本地（同一台机器）：}

\begin{itemize}
  \item 打开 \texttt{http://127.0.0.1:18789/}。
  \item 如果要求认证，运行 \texttt{openclaw dashboard} 并使用带令牌的链接（\texttt{?token=...}）。
  \item 令牌与 \texttt{gateway.auth.token}（或 \texttt{OPENCLAW\_GATEWAY\_TOKEN}）的值相同，UI 在首次加载后会存储它。
\end{itemize}


\textbf{非本地环境：}

\begin{itemize}
  \item \textbf{Tailscale Serve}（推荐）：保持绑定 loopback，运行 \texttt{openclaw gateway --tailscale serve}，打开 \texttt{https:///}。如果 \texttt{gateway.auth.allowTailscale} 为 \texttt{true}，身份标头满足认证要求（无需令牌）。
  \item \textbf{Tailnet 绑定}：运行 \texttt{openclaw gateway --bind tailnet --token ""}，打开 \texttt{http://:18789/}，在仪表板设置中粘贴令牌。
  \item \textbf{SSH 隧道}：\texttt{ssh -N -L 18789:127.0.0.1:18789 user@host}，然后从 \texttt{openclaw dashboard} 打开 \texttt{http://127.0.0.1:18789/?token=...}。
\end{itemize}


参阅\href{/web/dashboard}{仪表板}和 \href{/web}{Web 界面}了解绑定模式和认证详情。

\subsubsection{我需要什么运行时}


Node \textbf{>= 22} 是必需的。推荐使用 \texttt{pnpm}。\textbf{不推荐}使用 Bun 运行 Gateway 网关。

\subsubsection{能在 Raspberry Pi 上运行吗}


可以。Gateway 网关是轻量级的——文档列出 \textbf{512MB-1GB RAM}、\textbf{1 核}和约 \textbf{500MB} 磁盘空间足够个人使用，并指出 \textbf{Raspberry Pi 4 可以运行}。

如果你需要额外的余量（日志、媒体、其他服务），\textbf{推荐 2GB}，但这不是硬性最低要求。

提示：小型 Pi/VPS 可以托管 Gateway 网关，你可以在笔记本/手机上配对\textbf{节点}以获取本地屏幕/摄像头/画布或命令执行能力。参阅\href{/nodes}{节点}。

\subsubsection{Raspberry Pi 安装有什么建议}


简短回答：可以运行，但预期会有一些粗糙之处。

\begin{itemize}
  \item 使用 \textbf{64 位}操作系统并保持 Node >= 22。
  \item 优先选择\textbf{可编辑（git）安装}，以便查看日志和快速更新。
  \item 先不启用渠道/Skills，然后逐个添加。
  \item 如果遇到奇怪的二进制问题，通常是 \textbf{ARM 兼容性}问题。
\end{itemize}


文档：\href{/platforms/linux}{Linux}、\href{/install}{安装}。

\subsubsection{卡在 wake up my friend / 新手引导无法启动，怎么办}


该界面依赖于 Gateway 网关可达且已认证。TUI 也会在首次启动时自动发送"Wake up, my friend!"。如果你看到该行但\textbf{没有回复}且令牌保持为 0，说明智能体从未运行。

\begin{enumerate}
  \item 重启 Gateway 网关：
\end{enumerate}


\begin{lstlisting}[language=bash]
openclaw gateway restart
\end{lstlisting}


\begin{enumerate}
  \item 检查状态和认证：
\end{enumerate}


\begin{lstlisting}[language=bash]
openclaw status
openclaw models status
openclaw logs --follow
\end{lstlisting}


\begin{enumerate}
  \item 如果仍然挂起，运行：
\end{enumerate}


\begin{lstlisting}[language=bash]
openclaw doctor
\end{lstlisting}


如果 Gateway 网关在远程，确保隧道/Tailscale 连接正常，且 UI 指向正确的 Gateway 网关。参阅\href{/gateway/remote}{远程访问}。

\subsubsection{能否将我的设置迁移到新机器（Mac mini）而不重新进行新手引导}


可以。复制\textbf{状态目录}和\textbf{工作区}，然后运行一次 Doctor。只要你同时复制\textbf{两个}位置，就能保持你的机器人“完全一样”（记忆、会话历史、认证和渠道状态）：

\begin{enumerate}
  \item 在新机器上安装 OpenClaw。
  \item 从旧机器复制 \texttt{\$OPENCLAW\_STATE\_DIR}（默认：\texttt{\textasciitilde{}/.openclaw}）。
  \item 复制你的工作区（默认：\texttt{\textasciitilde{}/.openclaw/workspace}）。
  \item 运行 \texttt{openclaw doctor} 并重启 Gateway 网关服务。
\end{enumerate}


这会保留配置、认证配置文件、WhatsApp 凭据、会话和记忆。如果你处于远程模式，请记住 Gateway 网关主机拥有会话存储和工作区。

\textbf{重要：} 如果你只将工作区提交/推送到 GitHub，你只备份了\textbf{记忆 + 引导文件}，但\textbf{不包括}会话历史或认证。它们位于 \texttt{\textasciitilde{}/.openclaw/} 下（例如 \texttt{\textasciitilde{}/.openclaw/agents//sessions/}）。

相关：\href{/install/migrating}{迁移}、\href{/help/faq\#where-does-openclaw-store-its-data}{磁盘上的文件位置}、
\href{/concepts/agent-workspace}{智能体工作区}、\href{/gateway/doctor}{Doctor}、
\href{/gateway/remote}{远程模式}。

\subsubsection{在哪里查看最新版本的更新内容}


查看 GitHub 变更日志：
https://github.com/openclaw/openclaw/blob/main/CHANGELOG.md

最新条目在顶部。如果顶部部分标记为 \textbf{Unreleased}，则下一个带日期的部分是最新发布版本。条目按\textbf{亮点}、\textbf{变更}和\textbf{修复}分组（需要时还有文档/其他部分）。

\subsubsection{无法访问 docs.openclaw.ai（SSL 错误），怎么办}


一些 Comcast/Xfinity 连接通过 Xfinity Advanced Security 错误地拦截了 \texttt{docs.openclaw.ai}。禁用该功能或将 \texttt{docs.openclaw.ai} 加入白名单，然后重试。更多详情：\href{/help/troubleshooting\#docsopenclawai-shows-an-ssl-error-comcastxfinity}{故障排除}。
请帮助我们在此处报告以解除封锁：https://spa.xfinity.com/check\_url\_status。

如果仍然无法访问该网站，文档在 GitHub 上有镜像：
https://github.com/openclaw/openclaw/tree/main/docs

\subsubsection{stable 和 beta 有什么区别}


\textbf{Stable} 和 \textbf{beta} 是 \textbf{npm dist-tags}，不是独立的代码分支：

\begin{itemize}
  \item \texttt{latest} = stable
  \item \texttt{beta} = 用于测试的早期构建
\end{itemize}


我们将构建发布到 \textbf{beta}，测试后，一旦构建稳定，就会\textbf{将同一版本提升为 \texttt{latest}}。这就是为什么 beta 和 stable 可以指向\textbf{相同版本}。

查看变更：
https://github.com/openclaw/openclaw/blob/main/CHANGELOG.md

\subsubsection{如何安装 beta 版本，beta 和 dev 有什么区别}


\textbf{Beta} 是 npm dist-tag \texttt{beta}（可能与 \texttt{latest} 相同）。
\textbf{Dev} 是 \texttt{main} 的滚动头部（git）；发布时使用 npm dist-tag \texttt{dev}。

一行命令（macOS/Linux）：

\begin{lstlisting}[language=bash]
curl -fsSL --proto '=https' --tlsv1.2 https://openclaw.ai/install.sh | bash -s -- --beta
\end{lstlisting}


\begin{lstlisting}[language=bash]
curl -fsSL --proto '=https' --tlsv1.2 https://openclaw.ai/install.sh | bash -s -- --install-method git
\end{lstlisting}


Windows 安装程序（PowerShell）：
https://openclaw.ai/install.ps1

更多详情：\href{/install/development-channels}{开发渠道}和\href{/install/installer}{安装程序标志}。

\subsubsection{安装和新手引导通常需要多长时间}


大致指南：

\begin{itemize}
  \item \textbf{安装：} 2-5 分钟
  \item \textbf{新手引导：} 5-15 分钟，取决于配置多少渠道/模型
\end{itemize}


如果挂起，请参阅\href{/help/faq\#installer-stuck-how-do-i-get-more-feedback}{安装程序卡住}和\href{/help/faq\#im-stuck--whats-the-fastest-way-to-get-unstuck}{我卡住了}中的快速调试流程。

\subsubsection{如何试用最新代码}


两个选项：

\begin{enumerate}
  \item \textbf{Dev 渠道（git checkout）：}
\end{enumerate}


\begin{lstlisting}[language=bash]
openclaw update --channel dev
\end{lstlisting}


这会切换到 \texttt{main} 分支并从源码更新。

\begin{enumerate}
  \item \textbf{可编辑安装（从安装程序网站）：}
\end{enumerate}


\begin{lstlisting}[language=bash]
curl -fsSL https://openclaw.ai/install.sh | bash -s -- --install-method git
\end{lstlisting}


这会给你一个可编辑的本地仓库，然后通过 git 更新。

如果你更喜欢手动克隆，使用：

\begin{lstlisting}[language=bash]
git clone https://github.com/openclaw/openclaw.git
cd openclaw
pnpm install
pnpm build
\end{lstlisting}


文档：\href{/cli/update}{更新}、\href{/install/development-channels}{开发渠道}、
\href{/install}{安装}。

\subsubsection{安装程序卡住了？如何获取更多反馈}


使用\textbf{详细输出}重新运行安装程序：

\begin{lstlisting}[language=bash]
curl -fsSL https://openclaw.ai/install.sh | bash -s -- --verbose
\end{lstlisting}


带详细输出的 Beta 安装：

\begin{lstlisting}[language=bash]
curl -fsSL https://openclaw.ai/install.sh | bash -s -- --beta --verbose
\end{lstlisting}


可编辑（git）安装：

\begin{lstlisting}[language=bash]
curl -fsSL https://openclaw.ai/install.sh | bash -s -- --install-method git --verbose
\end{lstlisting}


更多选项：\href{/install/installer}{安装程序标志}。

\subsubsection{Windows 安装提示找不到 git 或无法识别 openclaw}


两个常见的 Windows 问题：

\textbf{1) npm error spawn git / git not found}

\begin{itemize}
  \item 安装 \textbf{Git for Windows} 并确保 \texttt{git} 在你的 PATH 中。
  \item 关闭并重新打开 PowerShell，然后重新运行安装程序。
\end{itemize}


\textbf{2) openclaw is not recognized（安装后）}

\begin{itemize}
  \item 你的 npm 全局 bin 文件夹不在 PATH 中。
  \item 检查路径：
\begin{lstlisting}
  npm config get prefix
\end{lstlisting}

  \item 确保 \texttt{\textbackslash\{\}\textbackslash\{\}bin} 在 PATH 中（在大多数系统上是 \texttt{\%AppData\%\textbackslash\{\}\textbackslash\{\}npm}）。
  \item 更新 PATH 后关闭并重新打开 PowerShell。
\end{itemize}


如果你想要最顺畅的 Windows 设置，请使用 \textbf{WSL2} 而不是原生 Windows。
文档：\href{/platforms/windows}{Windows}。

\subsubsection{文档没有解答我的问题——如何获得更好的答案}


使用\textbf{可编辑（git）安装}，这样你在本地拥有完整的源码和文档，然后从该文件夹向你的机器人（或 Claude/Codex）提问，这样它可以读取仓库并精确回答。

\begin{lstlisting}[language=bash]
curl -fsSL https://openclaw.ai/install.sh | bash -s -- --install-method git
\end{lstlisting}


更多详情：\href{/install}{安装}和\href{/install/installer}{安装程序标志}。

\subsubsection{如何在 Linux 上安装 OpenClaw}


简短回答：按照 Linux 指南操作，然后运行新手引导向导。

\begin{itemize}
  \item Linux 快速路径 + 服务安装：\href{/platforms/linux}{Linux}。
  \item 完整指南：\href{/start/getting-started}{入门}。
  \item 安装和更新：\href{/install/updating}{安装与更新}。
\end{itemize}


\subsubsection{如何在 VPS 上安装 OpenClaw}


任何 Linux VPS 都可以。在服务器上安装，然后使用 SSH/Tailscale 访问 Gateway 网关。

指南：\href{/install/exe-dev}{exe.dev}、\href{/install/hetzner}{Hetzner}、\href{/install/fly}{Fly.io}。
远程访问：\href{/gateway/remote}{Gateway 网关远程}。

\subsubsection{云/VPS 安装指南在哪里}


我们维护了一个\textbf{托管中心}，涵盖常见提供商。选择一个并按指南操作：

\begin{itemize}
  \item \href{/vps}{VPS 托管}（所有提供商汇总）
  \item \href{/install/fly}{Fly.io}
  \item \href{/install/hetzner}{Hetzner}
  \item \href{/install/exe-dev}{exe.dev}
\end{itemize}


在云端的工作方式：\textbf{Gateway 网关运行在服务器上}，你通过控制 UI（或 Tailscale/SSH）从笔记本/手机访问。你的状态 + 工作区位于服务器上，因此将主机视为数据来源并做好备份。

你可以将\textbf{节点}（Mac/iOS/Android/无头）配对到云端 Gateway 网关，以访问本地屏幕/摄像头/画布或在笔记本上执行命令，同时 Gateway 网关保持在云端。

中心：\href{/platforms}{平台}。远程访问：\href{/gateway/remote}{Gateway 网关远程}。
节点：\href{/nodes}{节点}、\href{/cli/nodes}{节点 CLI}。

\subsubsection{可以让 OpenClaw 自行更新吗}


简短回答：\textbf{可以，但不推荐}。更新流程可能重启 Gateway 网关（这会中断活跃会话），可能需要干净的 git checkout，并且可能提示确认。更安全的做法：作为运维人员从 shell 运行更新。

使用 CLI：

\begin{lstlisting}[language=bash]
openclaw update
openclaw update status
openclaw update --channel stable|beta|dev
openclaw update --tag <dist-tag|version>
openclaw update --no-restart
\end{lstlisting}


如果必须从智能体自动化：

\begin{lstlisting}[language=bash]
openclaw update --yes --no-restart
openclaw gateway restart
\end{lstlisting}


文档：\href{/cli/update}{更新}、\href{/install/updating}{更新指南}。

\subsubsection{新手引导向导具体做了什么}


\texttt{openclaw onboard} 是推荐的设置路径。在\textbf{本地模式}下，它引导你完成：

\begin{itemize}
  \item \textbf{模型/认证设置}（推荐使用 Anthropic \textbf{setup-token} 进行 Claude 订阅，支持 OpenAI Codex OAuth，API 密钥可选，支持 LM Studio 本地模型）
  \item \textbf{工作区}位置 + 引导文件
  \item \textbf{Gateway 网关设置}（绑定/端口/认证/tailscale）
  \item \textbf{渠道}（WhatsApp、Telegram、Discord、Mattermost（插件）、Signal、iMessage）
  \item \textbf{守护进程安装}（macOS 上的 LaunchAgent；Linux/WSL2 上的 systemd 用户单元）
  \item \textbf{健康检查}和\textbf{Skills}选择
\end{itemize}


如果你配置的模型未知或缺少认证，它还会发出警告。

\subsubsection{运行 OpenClaw 需要 Claude 或 OpenAI 订阅吗}


不需要。你可以使用 \textbf{API 密钥}（Anthropic/OpenAI/其他）或\textbf{纯本地模型}运行 OpenClaw，这样你的数据留在你的设备上。订阅（Claude Pro/Max 或 OpenAI Codex）是这些提供商的可选认证方式。

文档：\href{/providers/anthropic}{Anthropic}、\href{/providers/openai}{OpenAI}、
\href{/gateway/local-models}{本地模型}、\href{/concepts/models}{模型}。

\subsubsection{能否使用 Claude Max 订阅而不需要 API 密钥}


可以。你可以使用 \textbf{setup-token} 代替 API 密钥进行认证。这是订阅路径。

Claude Pro/Max 订阅\textbf{不包含 API 密钥}，因此这是订阅账户的正确方式。重要提示：你必须向 Anthropic 确认此用法是否符合其订阅政策和条款。如果你想要最明确、受支持的方式，请使用 Anthropic API 密钥。

\subsubsection{Anthropic setup-token 认证如何工作}


\texttt{claude setup-token} 通过 Claude Code CLI 生成一个\textbf{令牌字符串}（在 Web 控制台中不可用）。你可以在\textbf{任何机器}上运行它。在向导中选择 \textbf{Anthropic token (paste setup-token)} 或使用 \texttt{openclaw models auth paste-token --provider anthropic} 粘贴。令牌作为 \textbf{anthropic} 提供商的认证配置文件存储，像 API 密钥一样使用（无自动刷新）。更多详情：\href{/concepts/oauth}{OAuth}。

\subsubsection{在哪里获取 Anthropic setup-token}


它\textbf{不在} Anthropic Console 中。setup-token 由 \textbf{Claude Code CLI} 在\textbf{任何机器}上生成：

\begin{lstlisting}[language=bash]
claude setup-token
\end{lstlisting}


复制它打印的令牌，然后在向导中选择 \textbf{Anthropic token (paste setup-token)}。如果你想在 Gateway 网关主机上运行，使用 \texttt{openclaw models auth setup-token --provider anthropic}。如果你在其他地方运行了 \texttt{claude setup-token}，在 Gateway 网关主机上使用 \texttt{openclaw models auth paste-token --provider anthropic} 粘贴。参阅 \href{/providers/anthropic}{Anthropic}。

\subsubsection{是否支持 Claude 订阅认证（Claude Pro/Max）}


是的——通过 \textbf{setup-token}。OpenClaw 不再复用 Claude Code CLI OAuth 令牌；请使用 setup-token 或 Anthropic API 密钥。在任何地方生成令牌并在 Gateway 网关主机上粘贴。参阅 \href{/providers/anthropic}{Anthropic} 和 \href{/concepts/oauth}{OAuth}。

注意：Claude 订阅访问受 Anthropic 条款约束。对于生产或多用户工作负载，API 密钥通常是更安全的选择。

\subsubsection{为什么我看到 HTTP 429 rate\_limit\_error（来自 Anthropic）}


这意味着你当前窗口的 \textbf{Anthropic 配额/速率限制}已耗尽。如果你使用 \textbf{Claude 订阅}（setup-token 或 Claude Code OAuth），请等待窗口重置或升级你的计划。如果你使用 \textbf{Anthropic API 密钥}，请在 Anthropic Console 中检查使用量/计费并根据需要提高限制。

提示：设置一个\textbf{备用模型}，这样 OpenClaw 在某个提供商被限速时仍能继续回复。
参阅\href{/cli/models}{模型}和 \href{/concepts/oauth}{OAuth}。

\subsubsection{支持 AWS Bedrock 吗}


是的——通过 pi-ai 的 \textbf{Amazon Bedrock (Converse)} 提供商进行\textbf{手动配置}。你必须在 Gateway 网关主机上提供 AWS 凭据/区域，并在模型配置中添加 Bedrock 提供商条目。参阅 \href{/providers/bedrock}{Amazon Bedrock} 和\href{/providers/models}{模型提供商}。如果你更喜欢托管密钥流程，在 Bedrock 前面使用兼容 OpenAI 的代理仍然是有效选项。

\subsubsection{Codex 认证如何工作}


OpenClaw 通过 OAuth（ChatGPT 登录）支持 \textbf{OpenAI Code (Codex)}。向导可以运行 OAuth 流程，并在适当时将默认模型设置为 \texttt{openai-codex/gpt-5.2}。参阅\href{/concepts/model-providers}{模型提供商}和\href{/start/wizard}{向导}。

\subsubsection{是否支持 OpenAI 订阅认证（Codex OAuth）}


是的。OpenClaw 完全支持 \textbf{OpenAI Code (Codex) 订阅 OAuth}。新手引导向导可以为你运行 OAuth 流程。

参阅 \href{/concepts/oauth}{OAuth}、\href{/concepts/model-providers}{模型提供商}和\href{/start/wizard}{向导}。

\subsubsection{如何设置 Gemini CLI OAuth}


Gemini CLI 使用\textbf{插件认证流程}，而不是 \texttt{openclaw.json} 中的 client id 或 secret。

步骤：

\begin{enumerate}
  \item 启用插件：\texttt{openclaw plugins enable google-gemini-cli-auth}
  \item 登录：\texttt{openclaw models auth login --provider google-gemini-cli --set-default}
\end{enumerate}


这会在 Gateway 网关主机上将 OAuth 令牌存储为认证配置文件。详情：\href{/concepts/model-providers}{模型提供商}。

\subsubsection{本地模型适合日常聊天吗}


通常不适合。OpenClaw 需要大上下文 + 强安全性；小显卡会截断且泄漏。如果必须使用，请在本地运行你能运行的\textbf{最大} MiniMax M2.1 版本（LM Studio），参阅 \href{/gateway/local-models}{/gateway/local-models}。较小/量化的模型会增加提示注入风险——参阅\href{/gateway/security}{安全}。

\subsubsection{如何将托管模型流量限制在特定区域}


选择区域固定的端点。OpenRouter 为 MiniMax、Kimi 和 GLM 提供美国托管选项；选择美国托管变体以保持数据在区域内。你仍然可以通过使用 \texttt{models.mode: "merge"} 在这些旁边列出 Anthropic/OpenAI，这样故障转移保持可用，同时尊重你选择的区域提供商。

\subsubsection{我必须购买 Mac Mini 才能安装吗}


不需要。OpenClaw 运行在 macOS 或 Linux 上（Windows 通过 WSL2）。Mac mini 是可选的——有些人买一台作为常开主机，但小型 VPS、家庭服务器或 Raspberry Pi 级别的设备也可以。

你只有在使用 \textbf{macOS 专用工具}时才需要 Mac。对于 iMessage，你可以将 Gateway 网关保持在 Linux 上，通过将 \texttt{channels.imessage.cliPath} 指向 SSH 包装器在任何 Mac 上运行 \texttt{imsg}。如果你需要其他 macOS 专用工具，在 Mac 上运行 Gateway 网关或配对一个 macOS 节点。

文档：\href{/channels/imessage}{iMessage}、\href{/nodes}{节点}、\href{/platforms/mac/remote}{Mac 远程模式}。

\subsubsection{iMessage 支持需要 Mac mini 吗}


你需要\textbf{某台登录了 Messages 的 macOS 设备}。它\textbf{不一定}是 Mac mini——任何 Mac 都可以。OpenClaw 的 iMessage 集成在 macOS 上运行（BlueBubbles 或 \texttt{imsg}），而 Gateway 网关可以在其他地方运行。

常见设置：

\begin{itemize}
  \item 在 Linux/VPS 上运行 Gateway 网关，将 \texttt{channels.imessage.cliPath} 指向在 Mac 上运行 \texttt{imsg} 的 SSH 包装器。
  \item 如果你想要最简单的单机设置，在 Mac 上运行所有组件。
\end{itemize}


文档：\href{/channels/imessage}{iMessage}、\href{/channels/bluebubbles}{BlueBubbles}、
\href{/platforms/mac/remote}{Mac 远程模式}。

\subsubsection{如果我买了 Mac mini 运行 OpenClaw，能连接到我的 MacBook Pro 吗}


可以。\textbf{Mac mini 可以运行 Gateway 网关}，你的 MacBook Pro 可以作为\textbf{节点}（伴随设备）连接。节点不运行 Gateway 网关——它们提供额外功能，如该设备上的屏幕/摄像头/画布和 \texttt{system.run}。

常见模式：

\begin{itemize}
  \item Gateway 网关在 Mac mini 上（常开）。
  \item MacBook Pro 运行 macOS 应用或节点主机并配对到 Gateway 网关。
  \item 使用 \texttt{openclaw nodes status} / \texttt{openclaw nodes list} 查看它。
\end{itemize}


文档：\href{/nodes}{节点}、\href{/cli/nodes}{节点 CLI}。

\subsubsection{可以使用 Bun 吗}


Bun \textbf{不推荐}。我们观察到运行时 bug，特别是在 WhatsApp 和 Telegram 方面。
使用 \textbf{Node} 以获得稳定的 Gateway 网关。

如果你仍想尝试 Bun，请在没有 WhatsApp/Telegram 的非生产 Gateway 网关上进行。

\subsubsection{Telegram：allowFrom 填什么}


\texttt{channels.telegram.allowFrom} 是\textbf{人类发送者的 Telegram 用户 ID}（数字，推荐）或 \texttt{@username}。它不是机器人用户名。

更安全的方式（无需第三方机器人）：

\begin{itemize}
  \item 给你的机器人发私信，然后运行 \texttt{openclaw logs --follow} 并读取 \texttt{from.id}。
\end{itemize}


官方 Bot API：

\begin{itemize}
  \item 给你的机器人发私信，然后调用 \texttt{https://api.telegram.org/bot/getUpdates} 并读取 \texttt{message.from.id}。
\end{itemize}


第三方（隐私性较低）：

\begin{itemize}
  \item 给 \texttt{@userinfobot} 或 \texttt{@getidsbot} 发私信。
\end{itemize}


参阅 \href{/channels/telegram\#access-control-dms--groups}{/channels/telegram}。

\subsubsection{多人能否使用同一个 WhatsApp 号码配合不同的 OpenClaw 实例}


可以，通过\textbf{多智能体路由}。将每个发送者的 WhatsApp \textbf{私信}（peer \texttt{kind: "dm"}，发送者 E.164 格式如 \texttt{+15551234567}）绑定到不同的 \texttt{agentId}，这样每个人获得自己的工作区和会话存储。回复仍然来自\textbf{同一个 WhatsApp 账户}，且私信访问控制（\texttt{channels.whatsapp.dmPolicy} / \texttt{channels.whatsapp.allowFrom}）对每个 WhatsApp 账户是全局的。参阅\href{/concepts/multi-agent}{多智能体路由}和 \href{/channels/whatsapp}{WhatsApp}。

\subsubsection{能否同时运行一个“快速聊天”智能体和一个“用 Opus 编程”的智能体}


可以。使用多智能体路由：为每个智能体设置自己的默认模型，然后将入站路由（提供商账户或特定对等方）绑定到每个智能体。示例配置位于\href{/concepts/multi-agent}{多智能体路由}。另参阅\href{/concepts/models}{模型}和\href{/gateway/configuration}{配置}。

\subsubsection{Homebrew 在 Linux 上可用吗}


可以。Homebrew 支持 Linux（Linuxbrew）。快速设置：

\begin{lstlisting}[language=bash]
/bin/bash -c "$(curl -fsSL https://raw.githubusercontent.com/Homebrew/install/HEAD/install.sh)"
echo 'eval "$(/home/linuxbrew/.linuxbrew/bin/brew shellenv)"' >> ~/.profile
eval "$(/home/linuxbrew/.linuxbrew/bin/brew shellenv)"
brew install <formula>
\end{lstlisting}


如果你通过 systemd 运行 OpenClaw，确保服务 PATH 包含 \texttt{/home/linuxbrew/.linuxbrew/bin}（或你的 brew 前缀），以便 \texttt{brew} 安装的工具在非登录 shell 中可解析。
最近的构建还会在 Linux systemd 服务上自动添加常见的用户 bin 目录（例如 \texttt{\textasciitilde{}/.local/bin}、\texttt{\textasciitilde{}/.npm-global/bin}、\texttt{\textasciitilde{}/.local/share/pnpm}、\texttt{\textasciitilde{}/.bun/bin}），并在设置时尊重 \texttt{PNPM\_HOME}、\texttt{NPM\_CONFIG\_PREFIX}、\texttt{BUN\_INSTALL}、\texttt{VOLTA\_HOME}、\texttt{ASDF\_DATA\_DIR}、\texttt{NVM\_DIR} 和 \texttt{FNM\_DIR}。

\subsubsection{可编辑（git）安装和 npm 安装有什么区别}


\begin{itemize}
  \item \textbf{可编辑（git）安装：} 完整源码 checkout，可编辑，最适合贡献者。你在本地运行构建并可以修补代码/文档。
  \item \textbf{npm 安装：} 全局 CLI 安装，无仓库，最适合“直接运行”。更新来自 npm dist-tags。
\end{itemize}


文档：\href{/start/getting-started}{入门}、\href{/install/updating}{更新}。

\subsubsection{之后可以在 npm 和 git 安装之间切换吗}


可以。安装另一种方式，然后运行 Doctor 使 Gateway 网关服务指向新的入口点。
这\textbf{不会删除你的数据}——它只改变 OpenClaw 代码的安装位置。你的状态
（\texttt{\textasciitilde{}/.openclaw}）和工作区（\texttt{\textasciitilde{}/.openclaw/workspace}）保持不变。

从 npm → git：

\begin{lstlisting}[language=bash]
git clone https://github.com/openclaw/openclaw.git
cd openclaw
pnpm install
pnpm build
openclaw doctor
openclaw gateway restart
\end{lstlisting}


从 git → npm：

\begin{lstlisting}[language=bash]
npm install -g openclaw@latest
openclaw doctor
openclaw gateway restart
\end{lstlisting}


Doctor 会检测 Gateway 网关服务入口点不匹配，并提供重写服务配置以匹配当前安装的选项（在自动化中使用 \texttt{--repair}）。

备份提示：参阅\href{/help/faq\#whats-the-recommended-backup-strategy}{备份策略}。

\subsubsection{应该在笔记本电脑还是 VPS 上运行 Gateway 网关简短回答：**如果你想要 24/7 可靠性，使用 VPS**。如果你想要最低摩擦且能接受休眠/重启，在本地运行。}


\textbf{笔记本（本地 Gateway 网关）}

\begin{itemize}
  \item \textbf{优点：} 无服务器成本，直接访问本地文件，实时浏览器窗口。
  \item \textbf{缺点：} 休眠/网络中断 = 断连，操作系统更新/重启会中断，必须保持唤醒。
\end{itemize}


\textbf{VPS / 云}

\begin{itemize}
  \item \textbf{优点：} 常开，网络稳定，无笔记本休眠问题，更容易保持运行。
  \item \textbf{缺点：} 通常无头运行（使用截图），仅远程文件访问，更新需要 SSH。
\end{itemize}


\textbf{OpenClaw 特定说明：} WhatsApp/Telegram/Slack/Mattermost（插件）/Discord 在 VPS 上都能正常工作。唯一的真正权衡是\textbf{无头浏览器}与可见窗口。参阅\href{/tools/browser}{浏览器}。

\textbf{推荐默认值：} 如果之前遇到过 Gateway 网关断连，使用 VPS。当你正在积极使用 Mac 并且需要本地文件访问或可见浏览器的 UI 自动化时，本地运行很好。

\subsubsection{在专用机器上运行 OpenClaw 有多重要}


不是必需的，但\textbf{推荐用于可靠性和隔离}。

\begin{itemize}
  \item \textbf{专用主机（VPS/Mac mini/Pi）：} 常开，更少的休眠/重启中断，更干净的权限，更容易保持运行。
  \item \textbf{共享的笔记本/台式机：} 完全适合测试和活跃使用，但当机器休眠或更新时预期会有暂停。
\end{itemize}


如果你想要两全其美，将 Gateway 网关保持在专用主机上，并将笔记本配对为\textbf{节点}以获取本地屏幕/摄像头/执行工具。参阅\href{/nodes}{节点}。
安全指南请阅读\href{/gateway/security}{安全}。

\subsubsection{VPS 的最低要求和推荐操作系统是什么}


OpenClaw 是轻量级的。对于基本的 Gateway 网关 + 一个聊天渠道：

\begin{itemize}
  \item \textbf{绝对最低：} 1 vCPU，1GB RAM，约 500MB 磁盘。
  \item \textbf{推荐：} 1-2 vCPU，2GB RAM 或更多以留有余量（日志、媒体、多渠道）。节点工具和浏览器自动化可能消耗较多资源。
\end{itemize}


操作系统：使用 \textbf{Ubuntu LTS}（或任何现代 Debian/Ubuntu）。Linux 安装路径在那里测试得最充分。

文档：\href{/platforms/linux}{Linux}、\href{/vps}{VPS 托管}。

\subsubsection{可以在虚拟机中运行 OpenClaw 吗？有什么要求}


可以。将虚拟机视为与 VPS 相同：它需要常开、可达，并有足够的 RAM 用于 Gateway 网关和你启用的任何渠道。

基准指南：

\begin{itemize}
  \item \textbf{绝对最低：} 1 vCPU，1GB RAM。
  \item \textbf{推荐：} 2GB RAM 或更多，如果你运行多个渠道、浏览器自动化或媒体工具。
  \item \textbf{操作系统：} Ubuntu LTS 或其他现代 Debian/Ubuntu。
\end{itemize}


如果你使用 Windows，\textbf{WSL2 是最简单的虚拟机式设置}，具有最佳的工具兼容性。参阅 \href{/platforms/windows}{Windows}、\href{/vps}{VPS 托管}。
如果你在虚拟机中运行 macOS，参阅 \href{/install/macos-vm}{macOS VM}。

\subsection{什么是 OpenClaw？}


\subsubsection{用一段话描述 OpenClaw}


OpenClaw 是一个运行在你自己设备上的个人 AI 助手。它在你已经使用的消息平台上回复（WhatsApp、Telegram、Slack、Mattermost（插件）、Discord、Google Chat、Signal、iMessage、WebChat），还可以在支持的平台上进行语音和实时 Canvas。\textbf{Gateway 网关} 是常开的控制平面；助手是产品。

\subsubsection{价值主张是什么}


OpenClaw 不是“只是一个 Claude 包装器”。它是一个\textbf{本地优先的控制平面}，让你在\textbf{自己的硬件}上运行强大的助手，可从你已经使用的聊天应用访问，具有有状态会话、记忆和工具——无需将工作流程的控制权交给托管 SaaS。

亮点：

\begin{itemize}
  \item \textbf{你的设备，你的数据：} 在任何你想要的地方运行 Gateway 网关（Mac、Linux、VPS），并将工作区 + 会话历史保持在本地。
  \item \textbf{真实渠道，而非 Web 沙箱：} WhatsApp/Telegram/Slack/Discord/Signal/iMessage/等，加上支持平台上的移动语音和 Canvas。
  \item \textbf{模型无关：} 使用 Anthropic、OpenAI、MiniMax、OpenRouter 等，支持按智能体路由和故障转移。
  \item \textbf{纯本地选项：} 运行本地模型，让\textbf{所有数据都保留在你的设备上}。
  \item \textbf{多智能体路由：} 按渠道、账户或任务分配不同的智能体，每个都有自己的工作区和默认值。
  \item \textbf{开源且可编辑：} 无供应商锁定地检查、扩展和自托管。
\end{itemize}


文档：\href{/gateway}{Gateway 网关}、\href{/channels}{渠道}、\href{/concepts/multi-agent}{多智能体}、
\href{/concepts/memory}{记忆}。

\subsubsection{刚设置好，应该先做什么}


好的入门项目：

\begin{itemize}
  \item 建一个网站（WordPress、Shopify 或简单的静态站点）。
  \item 做一个移动应用原型（大纲、界面、API 计划）。
  \item 整理文件和文件夹（清理、命名、打标签）。
  \item 连接 Gmail 并自动化摘要或跟进。
\end{itemize}


它可以处理大型任务，但最好将其拆分为多个阶段，并使用子智能体进行并行工作。

\subsubsection{OpenClaw 日常最常用的五个场景是什么}


日常收益通常包括：

\begin{itemize}
  \item \textbf{个人简报：} 收件箱、日历和你关心的新闻摘要。
  \item \textbf{研究和起草：} 快速研究、摘要以及邮件或文档的初稿。
  \item \textbf{提醒和跟进：} 定时任务或心跳驱动的提醒和检查清单。
  \item \textbf{浏览器自动化：} 填写表单、收集数据和重复性网页任务。
  \item \textbf{跨设备协调：} 从手机发送任务，让 Gateway 网关在服务器上运行，然后在聊天中获取结果。
\end{itemize}


\subsubsection{OpenClaw 能否帮助 SaaS 进行获客、外联、广告和博客}


可以用于\textbf{调研、筛选和起草}。它可以扫描网站、建立候选名单、总结潜在客户，并撰写外联或广告文案草稿。

对于\textbf{外联或广告投放}，请保持人工审核。避免垃圾邮件，遵守当地法律和平台政策，在发送之前审查所有内容。最安全的模式是让 OpenClaw 起草，由你批准。

文档：\href{/gateway/security}{安全}。

\subsubsection{相比 Claude Code，在 Web 开发方面有什么优势}


OpenClaw 是一个\textbf{个人助手}和协调层，不是 IDE 替代品。使用 Claude Code 或 Codex 在仓库中进行最快的直接编码循环。当你需要持久记忆、跨设备访问和工具编排时，使用 OpenClaw。

优势：

\begin{itemize}
  \item 跨会话的\textbf{持久记忆 + 工作区}
  \item \textbf{多平台访问}（WhatsApp、Telegram、TUI、WebChat）
  \item \textbf{工具编排}（浏览器、文件、调度、钩子）
  \item \textbf{常开 Gateway 网关}（在 VPS 上运行，从任何地方交互）
  \item 用于本地浏览器/屏幕/摄像头/执行的\textbf{节点}
\end{itemize}


展示：https://openclaw.ai/showcase

\subsection{Skills 与自动化}


\subsubsection{如何自定义 Skills 而不弄脏仓库}


使用托管覆盖而不是编辑仓库副本。将你的更改放在 \texttt{\textasciitilde{}/.openclaw/skills//SKILL.md}（或通过 \texttt{\textasciitilde{}/.openclaw/openclaw.json} 中的 \texttt{skills.load.extraDirs} 添加文件夹）。优先级是 \texttt{/skills} > \texttt{\textasciitilde{}/.openclaw/skills} > 内置，所以托管覆盖优先生效而不会修改 git。只有值得上游合并的编辑才应该放在仓库中并作为 PR 提交。

\subsubsection{可以从自定义文件夹加载 Skills 吗}


可以。通过 \texttt{\textasciitilde{}/.openclaw/openclaw.json} 中的 \texttt{skills.load.extraDirs} 添加额外目录（最低优先级）。默认优先级保持不变：\texttt{/skills} → \texttt{\textasciitilde{}/.openclaw/skills} → 内置 → \texttt{skills.load.extraDirs}。\texttt{clawhub} 默认安装到 \texttt{./skills}，OpenClaw 将其视为 \texttt{/skills}。

\subsubsection{如何为不同任务使用不同模型}


目前支持的模式有：

\begin{itemize}
  \item \textbf{定时任务}：隔离的任务可以为每个任务设置 \texttt{model} 覆盖。
  \item \textbf{子智能体}：将任务路由到具有不同默认模型的独立智能体。
  \item \textbf{按需切换}：使用 \texttt{/model} 随时切换当前会话模型。
\end{itemize}


参阅\href{/automation/cron-jobs}{定时任务}、\href{/concepts/multi-agent}{多智能体路由}和\href{/tools/slash-commands}{斜杠命令}。

\subsubsection{机器人在执行繁重工作时卡住了，如何卸载任务}


使用\textbf{子智能体}处理长时间或并行任务。子智能体在自己的会话中运行，返回摘要，并保持你的主聊天响应。

要求你的机器人“为这个任务生成一个子智能体”或使用 \texttt{/subagents}。
在聊天中使用 \texttt{/status} 查看 Gateway 网关当前正在做什么（以及是否忙碌）。

令牌提示：长任务和子智能体都消耗令牌。如果关注成本，通过 \texttt{agents.defaults.subagents.model} 为子智能体设置更便宜的模型。

文档：\href{/tools/subagents}{子智能体}。

\subsubsection{定时任务或提醒没有触发，应该检查什么}


定时任务在 Gateway 网关进程内运行。如果 Gateway 网关没有持续运行，计划任务将不会运行。

检查清单：

\begin{itemize}
  \item 确认 cron 已启用（\texttt{cron.enabled}）且未设置 \texttt{OPENCLAW\_SKIP\_CRON}。
  \item 检查 Gateway 网关是否 24/7 运行（无休眠/重启）。
  \item 验证任务的时区设置（\texttt{--tz} 与主机时区）。
\end{itemize}


调试：

\begin{lstlisting}[language=bash]
openclaw cron run <jobId> --force
openclaw cron runs --id <jobId> --limit 50
\end{lstlisting}


文档：\href{/automation/cron-jobs}{定时任务}、\href{/automation/cron-vs-heartbeat}{定时任务 vs 心跳}。

\subsubsection{如何在 Linux 上安装 Skills}


使用 \textbf{ClawHub}（CLI）或将 Skills 放入你的工作区。macOS Skills UI 在 Linux 上不可用。
浏览 Skills：https://clawhub.com。

安装 ClawHub CLI（选择一个包管理器）：

\begin{lstlisting}[language=bash]
npm i -g clawhub
\end{lstlisting}


\begin{lstlisting}[language=bash]
pnpm add -g clawhub
\end{lstlisting}


\subsubsection{OpenClaw 能否按计划或在后台持续运行任务}


可以。使用 Gateway 网关调度器：

\begin{itemize}
  \item \textbf{定时任务}用于计划或重复任务（跨重启持久化）。
  \item \textbf{心跳}用于“主会话”定期检查。
  \item \textbf{隔离任务}用于自主智能体发布摘要或投递到聊天。
\end{itemize}


文档：\href{/automation/cron-jobs}{定时任务}、\href{/automation/cron-vs-heartbeat}{定时任务 vs 心跳}、
\href{/gateway/heartbeat}{心跳}。

\textbf{能否从 Linux 运行仅限 Apple/macOS 的 Skills}

不能直接运行。macOS Skills 受 \texttt{metadata.openclaw.os} 和所需二进制文件限制，Skills 只有在 \textbf{Gateway 网关主机}上符合条件时才会出现在系统提示中。在 Linux 上，\texttt{darwin} 专用 Skills（如 \texttt{apple-notes}、\texttt{apple-reminders}、\texttt{things-mac}）不会加载，除非你覆盖限制。

你有三种支持的模式：

\textbf{方案 A - 在 Mac 上运行 Gateway 网关（最简单）。}
在 macOS 二进制文件所在的地方运行 Gateway 网关，然后从 Linux 通过\href{\#how-do-i-run-openclaw-in-remote-mode-client-connects-to-a-gateway-elsewhere}{远程模式}或 Tailscale 连接。Skills 正常加载，因为 Gateway 网关主机是 macOS。

\textbf{方案 B - 使用 macOS 节点（无需 SSH）。}
在 Linux 上运行 Gateway 网关，配对一个 macOS 节点（菜单栏应用），并在 Mac 上将\textbf{节点运行命令}设置为“始终询问”或“始终允许”。当所需二进制文件存在于节点上时，OpenClaw 可以将 macOS 专用 Skills 视为符合条件。智能体通过 \texttt{nodes} 工具运行这些 Skills。如果你选择“始终询问”，在提示中批准“始终允许”会将该命令添加到允许列表。

\textbf{方案 C - 通过 SSH 代理 macOS 二进制文件（高级）。}
保持 Gateway 网关在 Linux 上，但使所需的 CLI 二进制文件解析为在 Mac 上运行的 SSH 包装器。然后覆盖 Skills 以允许 Linux 使其保持符合条件。

\begin{enumerate}
  \item 为二进制文件创建 SSH 包装器（示例：\texttt{imsg}）：
\begin{lstlisting}[language=bash]
   #!/usr/bin/env bash
   set -euo pipefail
   exec ssh -T user@mac-host /opt/homebrew/bin/imsg "$@"
\end{lstlisting}

  \item 将包装器放在 Linux 主机的 \texttt{PATH} 上（例如 \texttt{\textasciitilde{}/bin/imsg}）。
  \item 覆盖 Skills 元数据（工作区或 \texttt{\textasciitilde{}/.openclaw/skills}）以允许 Linux：
\begin{lstlisting}
   ---
   name: imsg
   description: iMessage/SMS CLI for listing chats, history, watch, and sending.
   metadata: { "openclaw": { "os": ["darwin", "linux"], "requires": { "bins": ["imsg"] } } }
   ---
\end{lstlisting}

  \item 开始新会话以刷新 Skills 快照。
\end{enumerate}


对于 iMessage，你也可以将 \texttt{channels.imessage.cliPath} 指向 SSH 包装器（OpenClaw 只需要 stdio）。参阅 \href{/channels/imessage}{iMessage}。

\subsubsection{有 Notion 或 HeyGen 集成吗}


目前没有内置集成。

选项：

\begin{itemize}
  \item \textbf{自定义 Skills / 插件：} 最适合可靠的 API 访问（Notion/HeyGen 都有 API）。
  \item \textbf{浏览器自动化：} 无需编码但更慢且更脆弱。
\end{itemize}


如果你想按客户保留上下文（代理工作流），一个简单的模式是：

\begin{itemize}
  \item 每个客户一个 Notion 页面（上下文 + 偏好 + 当前工作）。
  \item 在会话开始时要求智能体获取该页面。
\end{itemize}


如果你想要原生集成，请提交功能请求或构建一个针对这些 API 的 Skills。

安装 Skills：

\begin{lstlisting}[language=bash]
clawhub install <skill-slug>
clawhub update --all
\end{lstlisting}


ClawHub 安装到当前目录下的 \texttt{./skills}（或回退到你配置的 OpenClaw 工作区）；OpenClaw 在下一个会话中将其视为 \texttt{/skills}。对于跨智能体共享的 Skills，将它们放在 \texttt{\textasciitilde{}/.openclaw/skills//SKILL.md}。某些 Skills 期望通过 Homebrew 安装二进制文件；在 Linux 上意味着 Linuxbrew（参阅上面的 Homebrew Linux 常见问题条目）。参阅\href{/tools/skills}{Skills}和 \href{/tools/clawhub}{ClawHub}。

\subsubsection{如何安装用于浏览器接管的 Chrome 扩展}


使用内置安装程序，然后在 Chrome 中加载未打包的扩展：

\begin{lstlisting}[language=bash]
openclaw browser extension install
openclaw browser extension path
\end{lstlisting}


然后 Chrome → \texttt{chrome://extensions} → 启用“开发者模式” → “加载已解压的扩展程序” → 选择该文件夹。

完整指南（包括远程 Gateway 网关 + 安全注意事项）：\href{/tools/chrome-extension}{Chrome 扩展}

如果 Gateway 网关运行在与 Chrome 同一台机器上（默认设置），你通常\textbf{不需要}额外配置。
如果 Gateway 网关运行在其他地方，在运行浏览器的机器上运行一个节点主机，以便 Gateway 网关可以代理浏览器操作。
你仍然需要在要控制的标签页上点击扩展按钮（它不会自动附加）。

\subsection{沙箱与记忆}


\subsubsection{有专门的沙箱文档吗}


有。参阅\href{/gateway/sandboxing}{沙箱}。对于 Docker 特定设置（完整 Gateway 网关在 Docker 中或沙箱镜像），参阅 \href{/install/docker}{Docker}。

\textbf{能否让私信保持私密，但群组用一个智能体公开沙箱隔离}

可以——如果你的私密流量是\textbf{私信}而公开流量是\textbf{群组}。

使用 \texttt{agents.defaults.sandbox.mode: "non-main"}，这样群组/频道会话（非主键）在 Docker 中运行，而主私信会话保持在主机上。然后通过 \texttt{tools.sandbox.tools} 限制沙箱会话中可用的工具。

设置指南 + 示例配置：\href{/channels/groups\#pattern-personal-dms-public-groups-single-agent}{群组：个人私信 + 公开群组}

关键配置参考：\href{/gateway/configuration\#agentsdefaultssandbox}{Gateway 网关配置}

\subsubsection{如何将主机文件夹绑定到沙箱中}


将 \texttt{agents.defaults.sandbox.docker.binds} 设置为 \texttt{["host:path:mode"]}（例如 \texttt{"/home/user/src:/src:ro"}）。全局 + 按智能体的绑定会合并；当 \texttt{scope: "shared"} 时按智能体的绑定会被忽略。对于敏感内容使用 \texttt{:ro}，并记住绑定会绕过沙箱文件系统隔离。参阅\href{/gateway/sandboxing\#custom-bind-mounts}{沙箱}和\href{/gateway/sandbox-vs-tool-policy-vs-elevated\#bind-mounts-security-quick-check}{沙箱 vs 工具策略 vs 提权}了解示例和安全注意事项。

\subsubsection{记忆是如何工作的}


OpenClaw 记忆就是智能体工作区中的 Markdown 文件：

\begin{itemize}
  \item 每日笔记在 \texttt{memory/YYYY-MM-DD.md}
  \item 精选的长期笔记在 \texttt{MEMORY.md}（仅限主/私密会话）
\end{itemize}


OpenClaw 还会运行\textbf{静默的预压缩记忆刷新}，以提醒模型在自动压缩之前写入持久笔记。这只在工作区可写时运行（只读沙箱会跳过）。参阅\href{/concepts/memory}{记忆}。

\subsubsection{记忆总是遗忘，如何让它持久保存}


要求机器人\textbf{将事实写入记忆}。长期笔记属于 \texttt{MEMORY.md}，短期上下文放入 \texttt{memory/YYYY-MM-DD.md}。

这仍然是我们正在改进的领域。提醒模型存储记忆会有帮助；它会知道如何操作。如果它持续遗忘，验证 Gateway 网关每次运行时是否使用相同的工作区。

文档：\href{/concepts/memory}{记忆}、\href{/concepts/agent-workspace}{智能体工作区}。

\subsubsection{语义记忆搜索需要 OpenAI API 密钥吗}


只有在使用 \textbf{OpenAI embeddings} 时才需要。Codex OAuth 覆盖 chat/completions 但\textbf{不}授予 embeddings 访问权限，因此\textbf{使用 Codex 登录（OAuth 或 Codex CLI 登录）}对语义记忆搜索没有帮助。OpenAI embeddings 仍然需要真正的 API 密钥（\texttt{OPENAI\_API\_KEY} 或 \texttt{models.providers.openai.apiKey}）。

如果你没有明确设置提供商，OpenClaw 会在能解析 API 密钥（认证配置文件、\texttt{models.providers.*.apiKey} 或环境变量）时自动选择提供商。如果 OpenAI 密钥可解析则优先使用 OpenAI，否则如果 Gemini 密钥可解析则使用 Gemini。如果两个密钥都不可用，记忆搜索保持禁用直到你配置它。如果你配置了本地模型路径且存在，OpenClaw 优先使用 \texttt{local}。

如果你更想保持本地运行，设置 \texttt{memorySearch.provider = "local"}（可选 \texttt{memorySearch.fallback = "none"}）。如果你想使用 Gemini embeddings，设置 \texttt{memorySearch.provider = "gemini"} 并提供 \texttt{GEMINI\_API\_KEY}（或 \texttt{memorySearch.remote.apiKey}）。我们支持 \textbf{OpenAI、Gemini 或本地} embedding 模型——参阅\href{/concepts/memory}{记忆}了解设置详情。

\subsubsection{记忆是否永久保留？有什么限制}


记忆文件保存在磁盘上，持久存在直到你删除它们。限制是你的存储空间，而不是模型。\textbf{会话上下文}仍然受模型上下文窗口限制，所以长对话可能会压缩或截断。这就是记忆搜索存在的原因——它只将相关部分拉回上下文。

文档：\href{/concepts/memory}{记忆}、\href{/concepts/context}{上下文}。

\subsection{磁盘上的文件位置}


\subsubsection{OpenClaw 使用的所有数据都保存在本地吗}


不是——\textbf{OpenClaw 的状态是本地的}，但\textbf{外部服务仍然会看到你发送给它们的内容}。

\begin{itemize}
  \item \textbf{默认本地：} 会话、记忆文件、配置和工作区位于 Gateway 网关主机上（\texttt{\textasciitilde{}/.openclaw} + 你的工作区目录）。
  \item \textbf{必然远程：} 你发送给模型提供商（Anthropic/OpenAI/等）的消息会发送到它们的 API，聊天平台（WhatsApp/Telegram/Slack/等）在它们的服务器上存储消息数据。
  \item \textbf{你控制范围：} 使用本地模型可以将提示保留在你的机器上，但渠道流量仍然通过渠道的服务器。
\end{itemize}


相关：\href{/concepts/agent-workspace}{智能体工作区}、\href{/concepts/memory}{记忆}。

\subsubsection{OpenClaw 将数据存储在哪里}


所有内容位于 \texttt{\$OPENCLAW\_STATE\_DIR}（默认：\texttt{\textasciitilde{}/.openclaw}）下：


\begin{table}[h]
\centering
\small
\begin{tabular}{|l|l|}
\hline
路径 & 用途 \\
\hline
\texttt{\$OPENCLAW\_STATE\_DIR/openclaw.json} & 主配置（JSON5） \\
\hline
\texttt{\$OPENCLAW\_STATE\_DIR/credentials/oauth.json} & 旧版 OAuth 导入（首次使用时复制到认证配置文件） \\
\hline
\texttt{\$OPENCLAW\_STATE\_DIR/agents//agent/auth-profiles.json} & 认证配置文件（OAuth + API 密钥） \\
\hline
\texttt{\$OPENCLAW\_STATE\_DIR/agents//agent/auth.json} & 运行时认证缓存（自动管理） \\
\hline
\texttt{\$OPENCLAW\_STATE\_DIR/credentials/} & 提供商状态（例如 \texttt{whatsapp//creds.json}） \\
\hline
\texttt{\$OPENCLAW\_STATE\_DIR/agents/} & 按智能体的状态（agentDir + 会话） \\
\hline
\texttt{\$OPENCLAW\_STATE\_DIR/agents//sessions/} & 对话历史和状态（按智能体） \\
\hline
\texttt{\$OPENCLAW\_STATE\_DIR/agents//sessions/sessions.json} & 会话元数据（按智能体） \\
\hline
\end{tabular}
\end{table}



旧版单智能体路径：\texttt{\textasciitilde{}/.openclaw/agent/*}（通过 \texttt{openclaw doctor} 迁移）。

你的\textbf{工作区}（AGENTS.md、记忆文件、Skills 等）是独立的，通过 \texttt{agents.defaults.workspace} 配置（默认：\texttt{\textasciitilde{}/.openclaw/workspace}）。

\subsubsection{AGENTS.md / SOUL.md / USER.md / MEMORY.md 应该放在哪里}


这些文件位于\textbf{智能体工作区}中，而不是 \texttt{\textasciitilde{}/.openclaw}。

\begin{itemize}
  \item \textbf{工作区（按智能体）}：\texttt{AGENTS.md}、\texttt{SOUL.md}、\texttt{IDENTITY.md}、\texttt{USER.md}、
\end{itemize}

\texttt{MEMORY.md}（或 \texttt{memory.md}）、\texttt{memory/YYYY-MM-DD.md}、可选的 \texttt{HEARTBEAT.md}。
\begin{itemize}
  \item \textbf{状态目录（\texttt{\textasciitilde{}/.openclaw}）}：配置、凭据、认证配置文件、会话、日志和共享 Skills（\texttt{\textasciitilde{}/.openclaw/skills}）。
\end{itemize}


默认工作区是 \texttt{\textasciitilde{}/.openclaw/workspace}，可通过以下方式配置：

\begin{lstlisting}[language=json]
{
  agents: { defaults: { workspace: "~/.openclaw/workspace" } },
}
\end{lstlisting}


如果机器人在重启后“忘记”了内容，确认 Gateway 网关每次启动时都使用相同的工作区（记住：远程模式使用 \textbf{Gateway 网关主机的}工作区，而不是你本地笔记本的）。

提示：如果你想要一个持久的行为或偏好，要求机器人\textbf{将其写入 AGENTS.md 或 MEMORY.md}，而不是依赖聊天历史。

参阅\href{/concepts/agent-workspace}{智能体工作区}和\href{/concepts/memory}{记忆}。

\subsubsection{推荐的备份策略是什么}


将你的\textbf{智能体工作区}放入一个\textbf{私有} git 仓库，并备份到某个私有位置（例如 GitHub 私有仓库）。这会捕获记忆 + AGENTS/SOUL/USER 文件，让你以后可以恢复助手的“思维”。

\textbf{不要}提交 \texttt{\textasciitilde{}/.openclaw} 下的任何内容（凭据、会话、令牌）。如果你需要完整恢复，将工作区和状态目录分别备份（参阅上面的迁移问题）。

文档：\href{/concepts/agent-workspace}{智能体工作区}。

\subsubsection{如何完全卸载 OpenClaw}


参阅专门指南：\href{/install/uninstall}{卸载}。

\subsubsection{智能体可以在工作区外工作吗}


可以。工作区是\textbf{默认 cwd} 和记忆锚点，不是硬沙箱。相对路径在工作区内解析，但绝对路径可以访问其他主机位置，除非启用了沙箱。如果你需要隔离，使用 \href{/gateway/sandboxing}{\texttt{agents.defaults.sandbox}} 或按智能体的沙箱设置。如果你希望某个仓库作为默认工作目录，将该智能体的 \texttt{workspace} 指向仓库根目录。OpenClaw 仓库只是源代码；除非你有意要让智能体在其中工作，否则保持工作区独立。

示例（仓库作为默认 cwd）：

\begin{lstlisting}[language=json]
{
  agents: {
    defaults: {
      workspace: "~/Projects/my-repo",
    },
  },
}
\end{lstlisting}


\subsubsection{我处于远程模式——会话存储在哪里}


会话状态归 \textbf{Gateway 网关主机}所有。如果你处于远程模式，你关心的会话存储在远程机器上，而不是你的本地笔记本上。参阅\href{/concepts/session}{会话管理}。

\subsection{配置基础}


\subsubsection{配置文件是什么格式？在哪里}


OpenClaw 从 \texttt{\$OPENCLAW\_CONFIG\_PATH}（默认：\texttt{\textasciitilde{}/.openclaw/openclaw.json}）读取可选的 \textbf{JSON5} 配置：

\begin{lstlisting}
$OPENCLAW_CONFIG_PATH
\end{lstlisting}


如果文件不存在，使用安全的默认值（包括默认工作区 \texttt{\textasciitilde{}/.openclaw/workspace}）。

\subsubsection{我设置了 gateway.bind: "lan"（或 "tailnet"），现在什么都监听不了 / UI 显示未授权}


非 local loopback 绑定\textbf{需要认证}。配置 \texttt{gateway.auth.mode} + \texttt{gateway.auth.token}（或使用 \texttt{OPENCLAW\_GATEWAY\_TOKEN}）。

\begin{lstlisting}[language=json]
{
  gateway: {
    bind: "lan",
    auth: {
      mode: "token",
      token: "replace-me",
    },
  },
}
\end{lstlisting}


注意：

\begin{itemize}
  \item \texttt{gateway.remote.token} 仅用于\textbf{远程 CLI 调用}；它不启用本地 Gateway 网关认证。
  \item 控制 UI 通过 \texttt{connect.params.auth.token}（存储在应用/UI 设置中）进行认证。避免将令牌放在 URL 中。
\end{itemize}


\subsubsection{为什么现在在 localhost 也需要令牌}


向导默认生成 Gateway 网关令牌（即使在 local loopback 上），因此\textbf{本地 WS 客户端必须认证}。这阻止了其他本地进程调用 Gateway 网关。在控制 UI 设置（或你的客户端配置）中粘贴令牌以连接。

如果你\textbf{确实}想要开放 local loopback，从配置中移除 \texttt{gateway.auth}。Doctor 可以随时为你生成令牌：\texttt{openclaw doctor --generate-gateway-token}。

\subsubsection{更改配置后需要重启吗}


Gateway 网关监视配置文件并支持热重载：

\begin{itemize}
  \item \texttt{gateway.reload.mode: "hybrid"}（默认）：安全更改热应用，关键更改重启
  \item 也支持 \texttt{hot}、\texttt{restart}、\texttt{off}
\end{itemize}


\subsubsection{如何启用网络搜索（和网页抓取）}


\texttt{web\_fetch} 无需 API 密钥即可工作。\texttt{web\_search} 需要 Brave Search API 密钥。\textbf{推荐：} 运行 \texttt{openclaw configure --section web} 将其存储在 \texttt{tools.web.search.apiKey} 中。环境变量替代方案：为 Gateway 网关进程设置 \texttt{BRAVE\_API\_KEY}。

\begin{lstlisting}[language=json]
{
  tools: {
    web: {
      search: {
        enabled: true,
        apiKey: "BRAVE_API_KEY_HERE",
        maxResults: 5,
      },
      fetch: {
        enabled: true,
      },
    },
  },
}
\end{lstlisting}


注意：

\begin{itemize}
  \item 如果你使用允许列表，添加 \texttt{web\_search}/\texttt{web\_fetch} 或 \texttt{group:web}。
  \item \texttt{web\_fetch} 默认启用（除非明确禁用）。
  \item 守护进程从 \texttt{\textasciitilde{}/.openclaw/.env}（或服务环境）读取环境变量。
\end{itemize}


文档：\href{/tools/web}{Web 工具}。

\subsubsection{config.apply 清空了我的配置，如何恢复和避免}


\texttt{config.apply} 替换\textbf{整个配置}。如果你发送部分对象，其他所有内容都会被移除。

恢复：

\begin{itemize}
  \item 从备份恢复（git 或复制的 \texttt{\textasciitilde{}/.openclaw/openclaw.json}）。
  \item 如果没有备份，重新运行 \texttt{openclaw doctor} 并重新配置渠道/模型。
  \item 如果这是意外情况，提交 bug 并附上你最后已知的配置或任何备份。
  \item 本地编码智能体通常可以从日志或历史中重建工作配置。
\end{itemize}


避免方法：

\begin{itemize}
  \item 对小更改使用 \texttt{openclaw config set}。
  \item 对交互式编辑使用 \texttt{openclaw configure}。
\end{itemize}


文档：\href{/cli/config}{Config}、\href{/cli/configure}{Configure}、\href{/gateway/doctor}{Doctor}。

\subsubsection{如何运行一个中心 Gateway 网关配合跨设备的专用工作节点}


常见模式是\textbf{一个 Gateway 网关}（例如 Raspberry Pi）加上\textbf{节点}和\textbf{智能体}：

\begin{itemize}
  \item \textbf{Gateway 网关（中心）：} 拥有渠道（Signal/WhatsApp）、路由和会话。
  \item \textbf{节点（设备）：} Mac/iOS/Android 作为外围设备连接，暴露本地工具（\texttt{system.run}、\texttt{canvas}、\texttt{camera}）。
  \item \textbf{智能体（工作者）：} 用于特殊角色的独立大脑/工作区（例如“Hetzner 运维”、“个人数据”）。
  \item \textbf{子智能体：} 需要并行处理时从主智能体生成后台工作。
  \item \textbf{TUI：} 连接到 Gateway 网关并切换智能体/会话。
\end{itemize}


文档：\href{/nodes}{节点}、\href{/gateway/remote}{远程访问}、\href{/concepts/multi-agent}{多智能体路由}、\href{/tools/subagents}{子智能体}、\href{/web/tui}{TUI}。

\subsubsection{OpenClaw 浏览器可以无头运行吗}


可以。这是一个配置选项：

\begin{lstlisting}[language=json]
{
  browser: { headless: true },
  agents: {
    defaults: {
      sandbox: { browser: { headless: true } },
    },
  },
}
\end{lstlisting}


默认为 \texttt{false}（有头）。无头模式在某些网站上更容易触发反机器人检测。参阅\href{/tools/browser}{浏览器}。

无头模式使用\textbf{相同的 Chromium 引擎}，适用于大多数自动化（表单、点击、抓取、登录）。主要区别：

\begin{itemize}
  \item 没有可见的浏览器窗口（如果需要视觉效果使用截图）。
  \item 某些网站在无头模式下对自动化更严格（验证码、反机器人）。例如，X/Twitter 经常阻止无头会话。
\end{itemize}


\subsubsection{如何使用 Brave 进行浏览器控制}


将 \texttt{browser.executablePath} 设置为你的 Brave 二进制文件（或任何基于 Chromium 的浏览器）并重启 Gateway 网关。
参阅\href{/tools/browser\#use-brave-or-another-chromium-based-browser}{浏览器}中的完整配置示例。

\subsection{远程 Gateway 网关与节点}


\subsubsection{命令如何在 Telegram、Gateway 网关和节点之间传播}


Telegram 消息由 \textbf{Gateway 网关} 处理。Gateway 网关运行智能体，只有在需要节点工具时才通过 \textbf{Gateway 网关 WebSocket} 调用节点：

Telegram → Gateway 网关 → 智能体 → \texttt{node.*} → 节点 → Gateway 网关 → Telegram

节点不会看到入站提供商流量；它们只接收节点 RPC 调用。

\subsubsection{如果 Gateway 网关托管在远程，我的智能体如何访问我的电脑}


简短回答：\textbf{将你的电脑配对为节点}。Gateway 网关运行在其他地方，但它可以通过 Gateway 网关 WebSocket 在你的本地机器上调用 \texttt{node.*} 工具（屏幕、摄像头、系统）。

典型设置：

\begin{enumerate}
  \item 在常开主机（VPS/家庭服务器）上运行 Gateway 网关。
  \item 将 Gateway 网关主机和你的电脑放在同一个 tailnet 上。
  \item 确保 Gateway 网关 WS 可达（tailnet 绑定或 SSH 隧道）。
  \item 在本地打开 macOS 应用并以\textbf{远程 over SSH} 模式连接（或直接 tailnet），使其可以注册为节点。
  \item 在 Gateway 网关上批准节点：
\begin{lstlisting}[language=bash]
   openclaw nodes pending
   openclaw nodes approve <requestId>
\end{lstlisting}

\end{enumerate}


不需要单独的 TCP 桥接；节点通过 Gateway 网关 WebSocket 连接。

安全提醒：配对 macOS 节点允许在该机器上执行 \texttt{system.run}。只配对你信任的设备，并查阅\href{/gateway/security}{安全}。

文档：\href{/nodes}{节点}、\href{/gateway/protocol}{Gateway 网关协议}、\href{/platforms/mac/remote}{macOS 远程模式}、\href{/gateway/security}{安全}。

\subsubsection{Tailscale 已连接但收不到回复，怎么办}


检查基础项：

\begin{itemize}
  \item Gateway 网关正在运行：\texttt{openclaw gateway status}
  \item Gateway 网关健康：\texttt{openclaw status}
  \item 渠道健康：\texttt{openclaw channels status}
\end{itemize}


然后验证认证和路由：

\begin{itemize}
  \item 如果你使用 Tailscale Serve，确保 \texttt{gateway.auth.allowTailscale} 设置正确。
  \item 如果你通过 SSH 隧道连接，确认本地隧道已启动并指向正确端口。
  \item 确认你的允许列表（私信或群组）包含你的账户。
\end{itemize}


文档：\href{/gateway/tailscale}{Tailscale}、\href{/gateway/remote}{远程访问}、\href{/channels}{渠道}。

\subsubsection{两个 OpenClaw 实例（本地 + VPS）可以互相通信吗}


可以。没有内置的“机器人对机器人”桥接，但你可以通过几种可靠的方式实现：

\textbf{最简单：} 使用两个机器人都能访问的普通聊天渠道（Telegram/Slack/WhatsApp）。让机器人 A 给机器人 B 发消息，然后让机器人 B 正常回复。

\textbf{CLI 桥接（通用）：} 运行一个脚本调用另一个 Gateway 网关，使用 \texttt{openclaw agent --message ... --deliver}，定向到另一个机器人监听的聊天。如果一个机器人在远程 VPS 上，通过 SSH/Tailscale 将你的 CLI 指向该远程 Gateway 网关（参阅\href{/gateway/remote}{远程访问}）。

示例模式（从能到达目标 Gateway 网关的机器上运行）：

\begin{lstlisting}[language=bash]
openclaw agent --message "Hello from local bot" --deliver --channel telegram --reply-to <chat-id>
\end{lstlisting}


提示：添加护栏防止两个机器人无限循环（仅提及、渠道允许列表或“不回复机器人消息”规则）。

文档：\href{/gateway/remote}{远程访问}、\href{/cli/agent}{Agent CLI}、\href{/tools/agent-send}{Agent send}。

\subsubsection{多个智能体需要独立的 VPS 吗}


不需要。一个 Gateway 网关可以托管多个智能体，每个都有自己的工作区、模型默认值和路由。这是正常设置，比每个智能体一个 VPS 便宜且简单得多。

只有在需要硬隔离（安全边界）或非常不同的配置（你不想共享）时才使用独立的 VPS。否则保持一个 Gateway 网关并使用多个智能体或子智能体。

\subsubsection{在个人笔记本电脑上使用节点而不是从 VPS SSH 有什么好处}


有——节点是从远程 Gateway 网关到达你笔记本的首选方式，它们解锁的不仅仅是 shell 访问。Gateway 网关运行在 macOS/Linux（Windows 通过 WSL2）上且是轻量级的（小型 VPS 或 Raspberry Pi 级别的设备就够用；4 GB RAM 足够），所以常见设置是一个常开主机加上你的笔记本作为节点。

\begin{itemize}
  \item \textbf{无需入站 SSH。} 节点向 Gateway 网关 WebSocket 发起出站连接并使用设备配对。
  \item \textbf{更安全的执行控制。} \texttt{system.run} 受该笔记本上节点允许列表/审批的限制。
  \item \textbf{更多设备工具。} 节点除了 \texttt{system.run} 还暴露 \texttt{canvas}、\texttt{camera} 和 \texttt{screen}。
  \item \textbf{本地浏览器自动化。} 将 Gateway 网关保持在 VPS 上，但在本地运行 Chrome 并通过 Chrome 扩展 + 笔记本上的节点主机中继控制。
\end{itemize}


SSH 对临时 shell 访问很好，但节点对于持续的智能体工作流和设备自动化更简单。

文档：\href{/nodes}{节点}、\href{/cli/nodes}{节点 CLI}、\href{/tools/chrome-extension}{Chrome 扩展}。

\subsubsection{应该在第二台笔记本上安装还是只添加一个节点}


如果你只需要第二台笔记本上的\textbf{本地工具}（屏幕/摄像头/执行），将其添加为\textbf{节点}。这保持单一 Gateway 网关并避免重复配置。本地节点工具目前仅限 macOS，但我们计划扩展到其他操作系统。

只有在需要\textbf{硬隔离}或两个完全独立的机器人时才安装第二个 Gateway 网关。

文档：\href{/nodes}{节点}、\href{/cli/nodes}{节点 CLI}、\href{/gateway/multiple-gateways}{多 Gateway 网关}。

\subsubsection{节点会运行 Gateway 网关服务吗}


不会。每台主机上应该只运行\textbf{一个 Gateway 网关}，除非你有意运行隔离的配置文件（参阅\href{/gateway/multiple-gateways}{多 Gateway 网关}）。节点是连接到 Gateway 网关的外围设备（iOS/Android 节点，或 macOS 菜单栏应用的“节点模式”）。对于无头节点主机和 CLI 控制，参阅\href{/cli/node}{节点主机 CLI}。

\texttt{gateway}、\texttt{discovery} 和 \texttt{canvasHost} 的更改需要完全重启。

\subsubsection{有 API / RPC 方式来应用配置吗}


有。\texttt{config.apply} 验证 + 写入完整配置，并在操作过程中重启 Gateway 网关。

\subsubsection{首次安装的最小“合理”配置是什么}


\begin{lstlisting}[language=json]
{
  agents: { defaults: { workspace: "~/.openclaw/workspace" } },
  channels: { whatsapp: { allowFrom: ["+15555550123"] } },
}
\end{lstlisting}


这设置了你的工作区并限制谁可以触发机器人。

\subsubsection{如何在 VPS 上设置 Tailscale 并从 Mac 连接}


最简步骤：

\begin{enumerate}
  \item \textbf{在 VPS 上安装并登录}
\begin{lstlisting}[language=bash]
   curl -fsSL https://tailscale.com/install.sh | sh
   sudo tailscale up
\end{lstlisting}

  \item \textbf{在 Mac 上安装并登录}
\end{enumerate}

\begin{itemize}
  \item 使用 Tailscale 应用并登录到同一个 tailnet。
\end{itemize}

\begin{enumerate}
  \item \textbf{启用 MagicDNS（推荐）}
\end{enumerate}

\begin{itemize}
  \item 在 Tailscale 管理控制台中启用 MagicDNS，这样 VPS 有一个稳定的名称。
\end{itemize}

\begin{enumerate}
  \item \textbf{使用 tailnet 主机名}
\end{enumerate}

\begin{itemize}
  \item SSH：\texttt{ssh user@your-vps.tailnet-xxxx.ts.net}
  \item Gateway 网关 WS：\texttt{ws://your-vps.tailnet-xxxx.ts.net:18789}
\end{itemize}


如果你想要无 SSH 的控制 UI，在 VPS 上使用 Tailscale Serve：

\begin{lstlisting}[language=bash]
openclaw gateway --tailscale serve
\end{lstlisting}


这保持 Gateway 网关绑定到 local loopback 并通过 Tailscale 暴露 HTTPS。参阅 \href{/gateway/tailscale}{Tailscale}。

\subsubsection{如何将 Mac 节点连接到远程 Gateway 网关（Tailscale Serve）}


Serve 暴露 \textbf{Gateway 网关控制 UI + WS}。节点通过同一个 Gateway 网关 WS 端点连接。

推荐设置：

\begin{enumerate}
  \item \textbf{确保 VPS + Mac 在同一个 tailnet 上}。
  \item \textbf{使用 macOS 应用的远程模式}（SSH 目标可以是 tailnet 主机名）。应用会隧道 Gateway 网关端口并作为节点连接。
  \item \textbf{在 Gateway 网关上批准节点}：
\begin{lstlisting}[language=bash]
   openclaw nodes pending
   openclaw nodes approve <requestId>
\end{lstlisting}

\end{enumerate}


文档：\href{/gateway/protocol}{Gateway 网关协议}、\href{/gateway/discovery}{发现}、\href{/platforms/mac/remote}{macOS 远程模式}。

\subsection{环境变量和 .env 加载}


\subsubsection{OpenClaw 如何加载环境变量}


OpenClaw 从父进程（shell、launchd/systemd、CI 等）读取环境变量，并额外加载：

\begin{itemize}
  \item 当前工作目录下的 \texttt{.env}
  \item \texttt{\textasciitilde{}/.openclaw/.env}（即 \texttt{\$OPENCLAW\_STATE\_DIR/.env}）的全局回退 \texttt{.env}
\end{itemize}


两个 \texttt{.env} 文件都不会覆盖已有的环境变量。

你也可以在配置中定义内联环境变量（仅在进程环境中缺失时应用）：

\begin{lstlisting}[language=json]
{
  env: {
    OPENROUTER_API_KEY: "sk-or-...",
    vars: { GROQ_API_KEY: "gsk-..." },
  },
}
\end{lstlisting}


参阅 \href{/help/environment}{/environment} 了解优先级和来源详情。

\subsubsection{我通过服务启动了 Gateway 网关，但环境变量消失了，怎么办}


两个常见修复方法：

\begin{enumerate}
  \item 将缺失的密钥放在 \texttt{\textasciitilde{}/.openclaw/.env} 中，这样即使服务不继承你的 shell 环境也能被获取。
  \item 启用 shell 导入（可选的便利功能）：
\end{enumerate}


\begin{lstlisting}[language=json]
{
  env: {
    shellEnv: {
      enabled: true,
      timeoutMs: 15000,
    },
  },
}
\end{lstlisting}


这会运行你的登录 shell 并仅导入缺失的预期键名（从不覆盖）。环境变量等效项：
\texttt{OPENCLAW\_LOAD\_SHELL\_ENV=1}、\texttt{OPENCLAW\_SHELL\_ENV\_TIMEOUT\_MS=15000}。

\subsubsection{我设置了 COPILOT\_GITHUB\_TOKEN，但 models status 显示"Shell env: off"，为什么}


\texttt{openclaw models status} 报告的是 \textbf{shell 环境导入}是否启用。"Shell env: off"\textbf{不}意味着你的环境变量缺失——它只意味着 OpenClaw 不会自动加载你的登录 shell。

如果 Gateway 网关作为服务（launchd/systemd）运行，它不会继承你的 shell 环境。通过以下方式之一修复：

\begin{enumerate}
  \item 将令牌放在 \texttt{\textasciitilde{}/.openclaw/.env} 中：
\begin{lstlisting}
   COPILOT_GITHUB_TOKEN=...
\end{lstlisting}

  \item 或启用 shell 导入（\texttt{env.shellEnv.enabled: true}）。
  \item 或将其添加到配置的 \texttt{env} 块中（仅在缺失时应用）。
\end{enumerate}


然后重启 Gateway 网关并重新检查：

\begin{lstlisting}[language=bash]
openclaw models status
\end{lstlisting}


Copilot 令牌从 \texttt{COPILOT\_GITHUB\_TOKEN} 读取（也支持 \texttt{GH\_TOKEN} / \texttt{GITHUB\_TOKEN}）。
参阅 \href{/concepts/model-providers}{/concepts/model-providers} 和 \href{/help/environment}{/environment}。

\subsection{会话与多聊天}


\subsubsection{如何开始一个新对话}


发送 \texttt{/new} 或 \texttt{/reset} 作为独立消息。参阅\href{/concepts/session}{会话管理}。

\subsubsection{如果我从不发送 /new，会话会自动重置吗}


会。会话在 \texttt{session.idleMinutes}（默认 \textbf{60}）后过期。\textbf{下一条}消息会为该聊天键开始一个新的会话 ID。这不会删除记录——只是开始一个新会话。

\begin{lstlisting}[language=json]
{
  session: {
    idleMinutes: 240,
  },
}
\end{lstlisting}


\subsubsection{能否创建一个 OpenClaw 实例团队——一个 CEO 和多个智能体}


可以，通过\textbf{多智能体路由}和\textbf{子智能体}。你可以创建一个协调器智能体和多个工作者智能体，每个都有自己的工作区和模型。

话虽如此，最好将其视为一个\textbf{有趣的实验}。它消耗大量令牌，通常不如使用一个机器人配合不同会话的效率高。我们设想的典型模型是一个你与之对话的机器人，用不同的会话进行并行工作。该机器人也可以在需要时生成子智能体。

文档：\href{/concepts/multi-agent}{多智能体路由}、\href{/tools/subagents}{子智能体}、\href{/cli/agents}{智能体 CLI}。

\subsubsection{为什么上下文在任务中途被截断了？如何防止}


会话上下文受模型窗口限制。长对话、大量工具输出或许多文件可能触发压缩或截断。

有帮助的做法：

\begin{itemize}
  \item 要求机器人总结当前状态并写入文件。
  \item 在长任务之前使用 \texttt{/compact}，切换话题时使用 \texttt{/new}。
  \item 将重要上下文保存在工作区中，要求机器人读取。
  \item 对长时间或并行工作使用子智能体，这样主聊天保持较小。
  \item 如果这种情况经常发生，选择具有更大上下文窗口的模型。
\end{itemize}


\subsubsection{如何完全重置 OpenClaw 但保留安装}


使用重置命令：

\begin{lstlisting}[language=bash]
openclaw reset
\end{lstlisting}


非交互式完整重置：

\begin{lstlisting}[language=bash]
openclaw reset --scope full --yes --non-interactive
\end{lstlisting}


然后重新运行新手引导：

\begin{lstlisting}[language=bash]
openclaw onboard --install-daemon
\end{lstlisting}


注意：

\begin{itemize}
  \item 新手引导向导在看到现有配置时也提供\textbf{重置}选项。参阅\href{/start/wizard}{向导}。
  \item 如果你使用了配置文件（\texttt{--profile} / \texttt{OPENCLAW\_PROFILE}），重置每个状态目录（默认为 \texttt{\textasciitilde{}/.openclaw-}）。
  \item 开发重置：\texttt{openclaw gateway --dev --reset}（仅限开发；清除开发配置 + 凭据 + 会话 + 工作区）。
\end{itemize}


\subsubsection{我遇到了 context too large 错误——如何重置或压缩}


使用以下方式之一：

\begin{itemize}
  \item \textbf{压缩}（保留对话但总结较早的轮次）：
\end{itemize}


\begin{lstlisting}
  /compact
\end{lstlisting}


或 \texttt{/compact } 来引导总结。

\begin{itemize}
  \item \textbf{重置}（为同一聊天键开始新的会话 ID）：
\begin{lstlisting}
  /new
  /reset
\end{lstlisting}

\end{itemize}


如果持续出现：

\begin{itemize}
  \item 启用或调整\textbf{会话修剪}（\texttt{agents.defaults.contextPruning}）以裁剪旧的工具输出。
  \item 使用具有更大上下文窗口的模型。
\end{itemize}


文档：\href{/concepts/compaction}{压缩}、\href{/concepts/session-pruning}{会话修剪}、\href{/concepts/session}{会话管理}。

\subsubsection{为什么我看到 LLM request rejected: messages.N.content.X.tool\_use.input: Field required}


这是一个提供商验证错误：模型发出了一个没有必需 \texttt{input} 的 \texttt{tool\_use} 块。通常意味着会话历史已过时或损坏（通常在长线程或工具/模式变更后发生）。

修复：使用 \texttt{/new}（独立消息）开始新会话。

\subsubsection{为什么每 30 分钟收到一次心跳消息}


心跳默认每 \textbf{30 分钟}运行一次。调整或禁用：

\begin{lstlisting}[language=json]
{
  agents: {
    defaults: {
      heartbeat: {
        every: "2h", // 或 "0m" 禁用
      },
    },
  },
}
\end{lstlisting}


如果 \texttt{HEARTBEAT.md} 存在但实际上为空（只有空行和 markdown 标题如 \texttt{\# Heading}），OpenClaw 会跳过心跳运行以节省 API 调用。如果文件不存在，心跳仍然运行，由模型决定做什么。

按智能体覆盖使用 \texttt{agents.list[].heartbeat}。文档：\href{/gateway/heartbeat}{心跳}。

\subsubsection{需要在 WhatsApp 群组中添加“机器人账号”吗}


不需要。OpenClaw 运行在\textbf{你自己的账户}上，所以如果你在群组中，OpenClaw 就能看到它。
默认情况下，群组回复被阻止，直到你允许发送者（\texttt{groupPolicy: "allowlist"}）。

如果你只想\textbf{你自己}能触发群组回复：

\begin{lstlisting}[language=json]
{
  channels: {
    whatsapp: {
      groupPolicy: "allowlist",
      groupAllowFrom: ["+15551234567"],
    },
  },
}
\end{lstlisting}


\subsubsection{如何获取 WhatsApp 群组的 JID}


方法 1（最快）：跟踪日志并在群组中发送测试消息：

\begin{lstlisting}[language=bash]
openclaw logs --follow --json
\end{lstlisting}


查找以 \texttt{@g.us} 结尾的 \texttt{chatId}（或 \texttt{from}），如：
\texttt{1234567890-1234567890@g.us}。

方法 2（如果已配置/加入允许列表）：从配置中列出群组：

\begin{lstlisting}[language=bash]
openclaw directory groups list --channel whatsapp
\end{lstlisting}


文档：\href{/channels/whatsapp}{WhatsApp}、\href{/cli/directory}{目录}、\href{/cli/logs}{日志}。

\subsubsection{为什么 OpenClaw 不在群组中回复}


两个常见原因：

\begin{itemize}
  \item 提及限制已开启（默认）。你必须 @提及机器人（或匹配 \texttt{mentionPatterns}）。
  \item 你配置了 \texttt{channels.whatsapp.groups} 但没有 \texttt{"*"} 且该群组未加入允许列表。
\end{itemize}


参阅\href{/channels/groups}{群组}和\href{/channels/group-messages}{群组消息}。

\subsubsection{群组/线程与私聊共享上下文吗}


直接聊天默认折叠到主会话。群组/频道有自己的会话键，Telegram 话题 / Discord 线程是独立的会话。参阅\href{/channels/groups}{群组}和\href{/channels/group-messages}{群组消息}。

\subsubsection{可以创建多少个工作区和智能体}


没有硬性限制。几十个（甚至几百个）都没问题，但请注意：

\begin{itemize}
  \item \textbf{磁盘增长：} 会话 + 记录位于 \texttt{\textasciitilde{}/.openclaw/agents//sessions/} 下。
  \item \textbf{令牌成本：} 更多智能体意味着更多并发模型使用。
  \item \textbf{运维开销：} 按智能体的认证配置文件、工作区和渠道路由。
\end{itemize}


提示：

\begin{itemize}
  \item 每个智能体保持一个\textbf{活跃}工作区（\texttt{agents.defaults.workspace}）。
  \item 如果磁盘增长，修剪旧会话（删除 JSONL 或存储条目）。
  \item 使用 \texttt{openclaw doctor} 发现无用的工作区和配置文件不匹配。
\end{itemize}


\subsubsection{可以同时运行多个机器人或聊天（Slack）吗？应该如何设置}


可以。使用\textbf{多智能体路由}运行多个隔离的智能体，并按渠道/账户/对等方路由入站消息。Slack 作为渠道受支持，可以绑定到特定智能体。

浏览器访问功能强大，但不是“能做人类能做的一切”——反机器人、验证码和 MFA 仍然可以阻止自动化。为了最可靠的浏览器控制，在运行浏览器的机器上使用 Chrome 扩展中继（Gateway 网关可以在任何地方）。

最佳实践设置：

\begin{itemize}
  \item 常开 Gateway 网关主机（VPS/Mac mini）。
  \item 每个角色一个智能体（绑定）。
  \item Slack 渠道绑定到这些智能体。
  \item 需要时通过扩展中继（或节点）使用本地浏览器。
\end{itemize}


文档：\href{/concepts/multi-agent}{多智能体路由}、\href{/channels/slack}{Slack}、
\href{/tools/browser}{浏览器}、\href{/tools/chrome-extension}{Chrome 扩展}、\href{/nodes}{节点}。

\subsection{模型：默认值、选择、别名、切换}


\subsubsection{什么是“默认模型”}


OpenClaw 的默认模型是你设置的：

\begin{lstlisting}
agents.defaults.model.primary
\end{lstlisting}


模型以 \texttt{provider/model} 引用（示例：\texttt{anthropic/claude-opus-4-5}）。如果你省略提供商，OpenClaw 目前假设 \texttt{anthropic} 作为临时弃用回退——但你仍然应该\textbf{明确}设置 \texttt{provider/model}。

\subsubsection{推荐什么模型}


\textbf{推荐默认：} \texttt{anthropic/claude-opus-4-5}。
\textbf{好的替代：} \texttt{anthropic/claude-sonnet-4-5}。
\textbf{可靠（个性较少）：} \texttt{openai/gpt-5.2}——几乎和 Opus 一样好，只是个性较少。
\textbf{经济：} \texttt{zai/glm-4.7}。

MiniMax M2.1 有自己的文档：\href{/providers/minimax}{MiniMax} 和
\href{/gateway/local-models}{本地模型}。

经验法则：高风险工作使用你能负担的\textbf{最好的模型}，日常聊天或摘要使用更便宜的模型。你可以按智能体路由模型并使用子智能体并行处理长任务（每个子智能体消耗令牌）。参阅\href{/concepts/models}{模型}和\href{/tools/subagents}{子智能体}。

重要警告：较弱/过度量化的模型更容易受到提示注入和不安全行为的影响。参阅\href{/gateway/security}{安全}。

更多上下文：\href{/concepts/models}{模型}。

\subsubsection{可以使用自托管模型（llama.cpp、vLLM、Ollama）吗}


可以。如果你的本地服务器暴露了兼容 OpenAI 的 API，你可以将自定义提供商指向它。Ollama 直接支持，是最简单的路径。

安全说明：较小或大幅量化的模型更容易受到提示注入的影响。我们强烈建议对任何可以使用工具的机器人使用\textbf{大型模型}。如果你仍然想使用小模型，启用沙箱和严格的工具允许列表。

文档：\href{/providers/ollama}{Ollama}、\href{/gateway/local-models}{本地模型}、
\href{/concepts/model-providers}{模型提供商}、\href{/gateway/security}{安全}、
\href{/gateway/sandboxing}{沙箱}。

\subsubsection{如何在不清空配置的情况下切换模型}


使用\textbf{模型命令}或只编辑\textbf{模型}字段。避免完整配置替换。

安全选项：

\begin{itemize}
  \item 聊天中的 \texttt{/model}（快速，按会话）
  \item \texttt{openclaw models set ...}（只更新模型配置）
  \item \texttt{openclaw configure --section models}（交互式）
  \item 编辑 \texttt{\textasciitilde{}/.openclaw/openclaw.json} 中的 \texttt{agents.defaults.model}
\end{itemize}


避免使用部分对象执行 \texttt{config.apply}，除非你打算替换整个配置。如果你确实覆盖了配置，从备份恢复或重新运行 \texttt{openclaw doctor} 来修复。

文档：\href{/concepts/models}{模型}、\href{/cli/configure}{Configure}、\href{/cli/config}{Config}、\href{/gateway/doctor}{Doctor}。

\subsubsection{OpenClaw、Flawd 和 Krill 使用什么模型}


\begin{itemize}
  \item \textbf{OpenClaw + Flawd：} Anthropic Opus（\texttt{anthropic/claude-opus-4-5}）——参阅 \href{/providers/anthropic}{Anthropic}。
  \item \textbf{Krill：} MiniMax M2.1（\texttt{minimax/MiniMax-M2.1}）——参阅 \href{/providers/minimax}{MiniMax}。
\end{itemize}


\subsubsection{如何在运行中切换模型（无需重启）}


使用 \texttt{/model} 命令作为独立消息：

\begin{lstlisting}
/model sonnet
/model haiku
/model opus
/model gpt
/model gpt-mini
/model gemini
/model gemini-flash
\end{lstlisting}


你可以使用 \texttt{/model}、\texttt{/model list} 或 \texttt{/model status} 列出可用模型。

\texttt{/model}（和 \texttt{/model list}）显示紧凑的编号选择器。按编号选择：

\begin{lstlisting}
/model 3
\end{lstlisting}


你也可以为提供商强制指定特定的认证配置文件（按会话）：

\begin{lstlisting}
/model opus@anthropic:default
/model opus@anthropic:work
\end{lstlisting}


提示：\texttt{/model status} 显示哪个智能体是活跃的、正在使用哪个 \texttt{auth-profiles.json} 文件，以及接下来将尝试哪个认证配置文件。
它还显示配置的提供商端点（\texttt{baseUrl}）和 API 模式（\texttt{api}）（如果可用）。

\textbf{如何取消用 @profile 设置的配置文件固定}

重新运行 \texttt{/model} 但\textbf{不带} \texttt{@profile} 后缀：

\begin{lstlisting}
/model anthropic/claude-opus-4-5
\end{lstlisting}


如果你想返回默认值，从 \texttt{/model} 中选择（或发送 \texttt{/model }）。
使用 \texttt{/model status} 确认哪个认证配置文件是活跃的。

\subsubsection{能否日常任务用 GPT 5.2，编程用 Codex 5.2}


可以。设置一个为默认并按需切换：

\begin{itemize}
  \item \textbf{快速切换（按会话）：} 日常任务用 \texttt{/model gpt-5.2}，编程用 \texttt{/model gpt-5.2-codex}。
  \item \textbf{默认 + 切换：} 将 \texttt{agents.defaults.model.primary} 设置为 \texttt{openai-codex/gpt-5.2}，然后编程时切换到 \texttt{openai-codex/gpt-5.2-codex}（或反过来）。
  \item \textbf{子智能体：} 将编程任务路由到具有不同默认模型的子智能体。
\end{itemize}


参阅\href{/concepts/models}{模型}和\href{/tools/slash-commands}{斜杠命令}。

\subsubsection{为什么我看到"Model … is not allowed"然后没有回复}


如果设置了 \texttt{agents.defaults.models}，它成为 \texttt{/model} 和任何会话覆盖的\textbf{允许列表}。选择不在该列表中的模型会返回：

\begin{lstlisting}
Model "provider/model" is not allowed. Use /model to list available models.
\end{lstlisting}


该错误\textbf{代替}正常回复返回。修复：将模型添加到 \texttt{agents.defaults.models}，移除允许列表，或从 \texttt{/model list} 中选择一个模型。

\subsubsection{为什么我看到"Unknown model: minimax/MiniMax-M2.1"}


这意味着\textbf{提供商未配置}（未找到 MiniMax 提供商配置或认证配置文件），因此模型无法解析。此检测的修复在 \textbf{2026.1.12}（撰写本文时尚未发布）中。

修复清单：

\begin{enumerate}
  \item 升级到 \textbf{2026.1.12}（或从源码 \texttt{main} 运行），然后重启 Gateway 网关。
  \item 确保 MiniMax 已配置（向导或 JSON），或者 MiniMax API 密钥存在于环境/认证配置文件中以便提供商可以被注入。
  \item 使用精确的模型 ID（区分大小写）：\texttt{minimax/MiniMax-M2.1} 或 \texttt{minimax/MiniMax-M2.1-lightning}。
  \item 运行：
\begin{lstlisting}[language=bash]
   openclaw models list
\end{lstlisting}

\end{enumerate}

并从列表中选择（或在聊天中使用 \texttt{/model list}）。

参阅 \href{/providers/minimax}{MiniMax} 和\href{/concepts/models}{模型}。

\subsubsection{能否将 MiniMax 设为默认，复杂任务用 OpenAI}


可以。使用 \textbf{MiniMax 作为默认}，需要时\textbf{按会话}切换模型。故障转移用于\textbf{错误}，而非“困难任务”，所以使用 \texttt{/model} 或单独的智能体。

\textbf{方案 A：按会话切换}

\begin{lstlisting}[language=json]
{
  env: { MINIMAX_API_KEY: "sk-...", OPENAI_API_KEY: "sk-..." },
  agents: {
    defaults: {
      model: { primary: "minimax/MiniMax-M2.1" },
      models: {
        "minimax/MiniMax-M2.1": { alias: "minimax" },
        "openai/gpt-5.2": { alias: "gpt" },
      },
    },
  },
}
\end{lstlisting}


然后：

\begin{lstlisting}
/model gpt
\end{lstlisting}


\textbf{方案 B：分离智能体}

\begin{itemize}
  \item 智能体 A 默认：MiniMax
  \item 智能体 B 默认：OpenAI
  \item 按智能体路由或使用 \texttt{/agent} 切换
\end{itemize}


文档：\href{/concepts/models}{模型}、\href{/concepts/multi-agent}{多智能体路由}、\href{/providers/minimax}{MiniMax}、\href{/providers/openai}{OpenAI}。

\subsubsection{opus / sonnet / gpt 是内置快捷方式吗}


是的。OpenClaw 内置了一些默认简写（仅在模型存在于 \texttt{agents.defaults.models} 中时应用）：

\begin{itemize}
  \item \texttt{opus} → \texttt{anthropic/claude-opus-4-5}
  \item \texttt{sonnet} → \texttt{anthropic/claude-sonnet-4-5}
  \item \texttt{gpt} → \texttt{openai/gpt-5.2}
  \item \texttt{gpt-mini} → \texttt{openai/gpt-5-mini}
  \item \texttt{gemini} → \texttt{google/gemini-3-pro-preview}
  \item \texttt{gemini-flash} → \texttt{google/gemini-3-flash-preview}
\end{itemize}


如果你设置了同名的自定义别名，你的值优先。

\subsubsection{如何定义/覆盖模型快捷方式（别名）}


别名来自 \texttt{agents.defaults.models..alias}。示例：

\begin{lstlisting}[language=json]
{
  agents: {
    defaults: {
      model: { primary: "anthropic/claude-opus-4-5" },
      models: {
        "anthropic/claude-opus-4-5": { alias: "opus" },
        "anthropic/claude-sonnet-4-5": { alias: "sonnet" },
        "anthropic/claude-haiku-4-5": { alias: "haiku" },
      },
    },
  },
}
\end{lstlisting}


然后 \texttt{/model sonnet}（或支持时的 \texttt{/}）解析为该模型 ID。

\subsubsection{如何添加其他提供商（如 OpenRouter 或 Z.AI）的模型}


OpenRouter（按令牌付费；多种模型）：

\begin{lstlisting}[language=json]
{
  agents: {
    defaults: {
      model: { primary: "openrouter/anthropic/claude-sonnet-4-5" },
      models: { "openrouter/anthropic/claude-sonnet-4-5": {} },
    },
  },
  env: { OPENROUTER_API_KEY: "sk-or-..." },
}
\end{lstlisting}


Z.AI（GLM 模型）：

\begin{lstlisting}[language=json]
{
  agents: {
    defaults: {
      model: { primary: "zai/glm-4.7" },
      models: { "zai/glm-4.7": {} },
    },
  },
  env: { ZAI_API_KEY: "..." },
}
\end{lstlisting}


如果你引用了 provider/model 但缺少所需的提供商密钥，你会收到运行时认证错误（例如 \texttt{No API key found for provider "zai"}）。

\textbf{添加新智能体后提示 No API key found for provider}

这通常意味着\textbf{新智能体}的认证存储为空。认证是按智能体的，存储在：

\begin{lstlisting}
~/.openclaw/agents/<agentId>/agent/auth-profiles.json
\end{lstlisting}


修复选项：

\begin{itemize}
  \item 运行 \texttt{openclaw agents add } 并在向导中配置认证。
  \item 或从主智能体的 \texttt{agentDir} 复制 \texttt{auth-profiles.json} 到新智能体的 \texttt{agentDir}。
\end{itemize}


\textbf{不要}在智能体之间重用 \texttt{agentDir}；这会导致认证/会话冲突。

\subsection{模型故障转移与"All models failed"}


\subsubsection{故障转移是如何工作的}


故障转移分两个阶段：

\begin{enumerate}
  \item 同一提供商内的\textbf{认证配置文件轮换}。
  \item \textbf{模型回退}到 \texttt{agents.defaults.model.fallbacks} 中的下一个模型。
\end{enumerate}


冷却期适用于失败的配置文件（指数退避），因此 OpenClaw 即使在提供商被限速或临时失败时也能继续响应。

\subsubsection{这个错误是什么意思}


\begin{lstlisting}
No credentials found for profile "anthropic:default"
\end{lstlisting}


这意味着系统尝试使用认证配置文件 ID \texttt{anthropic:default}，但在预期的认证存储中找不到它的凭据。

\subsubsection{No credentials found for profile "anthropic:default" 的修复清单}


\begin{itemize}
  \item \textbf{确认认证配置文件的位置}（新路径 vs 旧路径）
  \item 当前：\texttt{\textasciitilde{}/.openclaw/agents//agent/auth-profiles.json}
  \item 旧版：\texttt{\textasciitilde{}/.openclaw/agent/*}（通过 \texttt{openclaw doctor} 迁移）
  \item \textbf{确认环境变量被 Gateway 网关加载}
  \item 如果你在 shell 中设置了 \texttt{ANTHROPIC\_API\_KEY} 但通过 systemd/launchd 运行 Gateway 网关，它可能不会继承。将其放在 \texttt{\textasciitilde{}/.openclaw/.env} 中或启用 \texttt{env.shellEnv}。
  \item \textbf{确保你编辑的是正确的智能体}
  \item 多智能体设置意味着可能有多个 \texttt{auth-profiles.json} 文件。
  \item \textbf{完整性检查模型/认证状态}
  \item 使用 \texttt{openclaw models status} 查看已配置的模型以及提供商是否已认证。
\end{itemize}


\textbf{No credentials found for profile "anthropic" 的修复清单}

这意味着运行固定到 Anthropic 认证配置文件，但 Gateway 网关在其认证存储中找不到它。

\begin{itemize}
  \item \textbf{使用 setup-token}
  \item 运行 \texttt{claude setup-token}，然后用 \texttt{openclaw models auth setup-token --provider anthropic} 粘贴。
  \item 如果令牌在另一台机器上创建，使用 \texttt{openclaw models auth paste-token --provider anthropic}。
  \item \textbf{如果你想使用 API 密钥}
  \item 在 \textbf{Gateway 网关主机}上将 \texttt{ANTHROPIC\_API\_KEY} 放入 \texttt{\textasciitilde{}/.openclaw/.env}。
  \item 清除任何强制缺失配置文件的固定顺序：
\begin{lstlisting}[language=bash]
    openclaw models auth order clear --provider anthropic
\end{lstlisting}

  \item \textbf{确认你在 Gateway 网关主机上运行命令}
  \item 在远程模式下，认证配置文件位于 Gateway 网关机器上，而不是你的笔记本上。
\end{itemize}


\subsubsection{为什么还尝试了 Google Gemini 并且失败了}


如果你的模型配置包含 Google Gemini 作为回退（或你切换到了 Gemini 简写），OpenClaw 会在模型回退期间尝试它。如果你没有配置 Google 凭据，你会看到 \texttt{No API key found for provider "google"}。

修复：提供 Google 认证，或从 \texttt{agents.defaults.model.fallbacks} / 别名中移除/避免 Google 模型，这样回退不会路由到那里。

\textbf{LLM request rejected message thinking signature required google antigravity}

原因：会话历史包含\textbf{没有签名的 thinking 块}（通常来自中止/部分流）。Google Antigravity 要求 thinking 块有签名。

修复：OpenClaw 现在为 Google Antigravity Claude 剥离未签名的 thinking 块。如果仍然出现，开始\textbf{新会话}或为该智能体设置 \texttt{/thinking off}。

\subsection{认证配置文件：概念和管理方式}


相关：\href{/concepts/oauth}{/concepts/oauth}（OAuth 流程、令牌存储、多账户模式）

\subsubsection{什么是认证配置文件}


认证配置文件是绑定到提供商的命名凭据记录（OAuth 或 API 密钥）。配置文件位于：

\begin{lstlisting}
~/.openclaw/agents/<agentId>/agent/auth-profiles.json
\end{lstlisting}


\subsubsection{典型的配置文件 ID 有哪些}


OpenClaw 使用提供商前缀的 ID，如：

\begin{itemize}
  \item \texttt{anthropic:default}（没有邮箱身份时常见）
  \item \texttt{anthropic:}（用于 OAuth 身份）
  \item 你自定义的 ID（例如 \texttt{anthropic:work}）
\end{itemize}


\subsubsection{可以控制首先尝试哪个认证配置文件吗}


可以。配置支持配置文件的可选元数据和按提供商的排序（\texttt{auth.order.}）。这\textbf{不}存储密钥；它将 ID 映射到 provider/mode 并设置轮换顺序。

如果某个配置文件处于短期\textbf{冷却}（速率限制/超时/认证失败）或较长的\textbf{禁用}状态（计费/额度不足），OpenClaw 可能会临时跳过它。要检查这一点，运行 \texttt{openclaw models status --json} 并查看 \texttt{auth.unusableProfiles}。调优：\texttt{auth.cooldowns.billingBackoffHours*}。

你也可以通过 CLI 设置\textbf{按智能体}的顺序覆盖（存储在该智能体的 \texttt{auth-profiles.json} 中）：

\begin{lstlisting}[language=bash]
# 默认为配置的默认智能体（省略 --agent）
openclaw models auth order get --provider anthropic

# 将轮换锁定到单个配置文件（只尝试这一个）
openclaw models auth order set --provider anthropic anthropic:default

# 或设置明确的顺序（提供商内回退）
openclaw models auth order set --provider anthropic anthropic:work anthropic:default

# 清除覆盖（回退到配置 auth.order / 轮换）
openclaw models auth order clear --provider anthropic
\end{lstlisting}


要针对特定智能体：

\begin{lstlisting}[language=bash]
openclaw models auth order set --provider anthropic --agent main anthropic:default
\end{lstlisting}


\subsubsection{OAuth 与 API 密钥：有什么区别}


OpenClaw 两者都支持：

\begin{itemize}
  \item \textbf{OAuth} 通常利用订阅访问（如适用）。
  \item \textbf{API 密钥} 使用按令牌付费的计费。
\end{itemize}


向导明确支持 Anthropic setup-token 和 OpenAI Codex OAuth，也可以为你存储 API 密钥。

\subsection{Gateway 网关：端口、“已在运行”和远程模式}


\subsubsection{Gateway 网关使用什么端口}


\texttt{gateway.port} 控制用于 WebSocket + HTTP（控制 UI、钩子等）的单个复用端口。

优先级：

\begin{lstlisting}
--port > OPENCLAW_GATEWAY_PORT > gateway.port > 默认 18789
\end{lstlisting}


\subsubsection{为什么 openclaw gateway status 显示 Runtime: running 但 RPC probe: failed}


因为"running"是 \textbf{supervisor} 的视角（launchd/systemd/schtasks）。RPC 探测是 CLI 实际连接到 Gateway 网关 WebSocket 并调用 \texttt{status}。

使用 \texttt{openclaw gateway status} 并关注这些行：

\begin{itemize}
  \item \texttt{Probe target:}（探测实际使用的 URL）
  \item \texttt{Listening:}（端口上实际绑定的内容）
  \item \texttt{Last gateway error:}（进程存活但端口未监听时的常见根因）
\end{itemize}


\subsubsection{为什么 openclaw gateway status 显示 Config (cli) 和 Config (service) 不同}


你正在编辑一个配置文件，而服务运行的是另一个（通常是 \texttt{--profile} / \texttt{OPENCLAW\_STATE\_DIR} 不匹配）。

修复：

\begin{lstlisting}[language=bash]
openclaw gateway install --force
\end{lstlisting}


从你希望服务使用的相同 \texttt{--profile} / 环境运行该命令。

\subsubsection{"another gateway instance is already listening"是什么意思}


OpenClaw 通过在启动时立即绑定 WebSocket 监听器来强制运行时锁（默认 \texttt{ws://127.0.0.1:18789}）。如果绑定因 \texttt{EADDRINUSE} 失败，它会抛出 \texttt{GatewayLockError} 表示另一个实例已在监听。

修复：停止另一个实例，释放端口，或使用 \texttt{openclaw gateway --port } 运行。

\subsubsection{如何以远程模式运行 OpenClaw（客户端连接到其他位置的 Gateway 网关）}


设置 \texttt{gateway.mode: "remote"} 并指向远程 WebSocket URL，可选带令牌/密码：

\begin{lstlisting}[language=json]
{
  gateway: {
    mode: "remote",
    remote: {
      url: "ws://gateway.tailnet:18789",
      token: "your-token",
      password: "your-password",
    },
  },
}
\end{lstlisting}


注意：

\begin{itemize}
  \item \texttt{openclaw gateway} 仅在 \texttt{gateway.mode} 为 \texttt{local} 时启动（或你传递覆盖标志）。
  \item macOS 应用监视配置文件，当这些值更改时实时切换模式。
\end{itemize}


\subsubsection{控制 UI 显示"unauthorized"或持续重连，怎么办}


你的 Gateway 网关运行时启用了认证（\texttt{gateway.auth.*}），但 UI 没有发送匹配的令牌/密码。

事实（来自代码）：

\begin{itemize}
  \item 控制 UI 将令牌存储在浏览器 localStorage 键 \texttt{openclaw.control.settings.v1} 中。
  \item UI 可以导入一次 \texttt{?token=...}（和/或 \texttt{?password=...}），然后从 URL 中剥离。
\end{itemize}


修复：

\begin{itemize}
  \item 最快：\texttt{openclaw dashboard}（打印 + 复制带令牌的链接，尝试打开；如果无头则显示 SSH 提示）。
  \item 如果你还没有令牌：\texttt{openclaw doctor --generate-gateway-token}。
  \item 如果是远程，先建隧道：\texttt{ssh -N -L 18789:127.0.0.1:18789 user@host} 然后打开 \texttt{http://127.0.0.1:18789/?token=...}。
  \item 在 Gateway 网关主机上设置 \texttt{gateway.auth.token}（或 \texttt{OPENCLAW\_GATEWAY\_TOKEN}）。
  \item 在控制 UI 设置中粘贴相同的令牌（或使用一次性 \texttt{?token=...} 链接刷新）。
  \item 仍然卡住？运行 \texttt{openclaw status --all} 并按\href{/gateway/troubleshooting}{故障排除}操作。参阅\href{/web/dashboard}{仪表板}了解认证详情。
\end{itemize}


\subsubsection{我设置了 gateway.bind: "tailnet" 但无法绑定 / 什么都没监听}


\texttt{tailnet} 绑定从你的网络接口中选择 Tailscale IP（100.64.0.0/10）。如果机器没有在 Tailscale 上（或接口已关闭），就没有可绑定的地址。

修复：

\begin{itemize}
  \item 在该主机上启动 Tailscale（使其拥有 100.x 地址），或
  \item 切换到 \texttt{gateway.bind: "loopback"} / \texttt{"lan"}。
\end{itemize}


注意：\texttt{tailnet} 是明确的。\texttt{auto} 优先 local loopback；当你想要仅 tailnet 绑定时使用 \texttt{gateway.bind: "tailnet"}。

\subsubsection{可以在同一主机上运行多个 Gateway 网关吗}


通常不需要——一个 Gateway 网关可以运行多个消息渠道和智能体。仅在需要冗余（例如救援机器人）或硬隔离时使用多个 Gateway 网关。

可以，但你必须隔离：

\begin{itemize}
  \item \texttt{OPENCLAW\_CONFIG\_PATH}（每实例配置）
  \item \texttt{OPENCLAW\_STATE\_DIR}（每实例状态）
  \item \texttt{agents.defaults.workspace}（工作区隔离）
  \item \texttt{gateway.port}（唯一端口）
\end{itemize}


快速设置（推荐）：

\begin{itemize}
  \item 每实例使用 \texttt{openclaw --profile  …}（自动创建 \texttt{\textasciitilde{}/.openclaw-}）。
  \item 在每个配置文件配置中设置唯一的 \texttt{gateway.port}（或手动运行时传 \texttt{--port}）。
  \item 安装每配置文件的服务：\texttt{openclaw --profile  gateway install}。
\end{itemize}


配置文件还会为服务名称添加后缀（\texttt{bot.molt.}；旧版 \texttt{com.openclaw.*}、\texttt{openclaw-gateway-.service}、\texttt{OpenClaw Gateway 网关 ()}）。
完整指南：\href{/gateway/multiple-gateways}{多 Gateway 网关}。

\subsubsection{"invalid handshake" / code 1008 是什么意思}


Gateway 网关是一个 \textbf{WebSocket 服务器}，它期望第一条消息是 \texttt{connect} 帧。如果收到其他内容，它会以 \textbf{code 1008}（策略违规）关闭连接。

常见原因：

\begin{itemize}
  \item 你在浏览器中打开了 \textbf{HTTP} URL（\texttt{http://...}）而不是 WS 客户端。
  \item 你使用了错误的端口或路径。
  \item 代理或隧道剥离了认证头或发送了非 Gateway 网关请求。
\end{itemize}


快速修复：

\begin{enumerate}
  \item 使用 WS URL：\texttt{ws://:18789}（或 \texttt{wss://...} 如果 HTTPS）。
  \item 不要在普通浏览器标签页中打开 WS 端口。
  \item 如果认证已启用，在 \texttt{connect} 帧中包含令牌/密码。
\end{enumerate}


如果你使用 CLI 或 TUI，URL 应该类似：

\begin{lstlisting}
openclaw tui --url ws://<host>:18789 --token <token>
\end{lstlisting}


协议详情：\href{/gateway/protocol}{Gateway 网关协议}。

\subsection{日志与调试}


\subsubsection{日志在哪里}


文件日志（结构化）：

\begin{lstlisting}
/tmp/openclaw/openclaw-YYYY-MM-DD.log
\end{lstlisting}


你可以通过 \texttt{logging.file} 设置稳定路径。文件日志级别由 \texttt{logging.level} 控制。控制台详细度由 \texttt{--verbose} 和 \texttt{logging.consoleLevel} 控制。

最快的日志跟踪：

\begin{lstlisting}[language=bash]
openclaw logs --follow
\end{lstlisting}


服务/supervisor 日志（当 Gateway 网关通过 launchd/systemd 运行时）：

\begin{itemize}
  \item macOS：\texttt{\$OPENCLAW\_STATE\_DIR/logs/gateway.log} 和 \texttt{gateway.err.log}（默认：\texttt{\textasciitilde{}/.openclaw/logs/...}；配置文件使用 \texttt{\textasciitilde{}/.openclaw-/logs/...}）
  \item Linux：\texttt{journalctl --user -u openclaw-gateway[-].service -n 200 --no-pager}
  \item Windows：\texttt{schtasks /Query /TN "OpenClaw Gateway 网关 ()" /V /FO LIST}
\end{itemize}


参阅\href{/gateway/troubleshooting\#log-locations}{故障排除}了解更多。

\subsubsection{如何启动/停止/重启 Gateway 网关服务}


使用 Gateway 网关辅助命令：

\begin{lstlisting}[language=bash]
openclaw gateway status
openclaw gateway restart
\end{lstlisting}


如果你手动运行 Gateway 网关，\texttt{openclaw gateway --force} 可以回收端口。参阅 \href{/gateway}{Gateway 网关}。

\subsubsection{我在 Windows 上关闭了终端——如何重启 OpenClaw}


有\textbf{两种 Windows 安装模式}：

\textbf{1) WSL2（推荐）：} Gateway 网关运行在 Linux 内部。

打开 PowerShell，进入 WSL，然后重启：

\begin{lstlisting}
wsl
openclaw gateway status
openclaw gateway restart
\end{lstlisting}


如果你从未安装服务，在前台启动：

\begin{lstlisting}[language=bash]
openclaw gateway run
\end{lstlisting}


\textbf{2) 原生 Windows（不推荐）：} Gateway 网关直接在 Windows 中运行。

打开 PowerShell 并运行：

\begin{lstlisting}
openclaw gateway status
openclaw gateway restart
\end{lstlisting}


如果你手动运行（无服务），使用：

\begin{lstlisting}
openclaw gateway run
\end{lstlisting}


文档：\href{/platforms/windows}{Windows (WSL2)}、\href{/gateway}{Gateway 网关服务运维手册}。

\subsubsection{Gateway 网关已启动但回复始终不到达，应该检查什么}


从快速健康扫描开始：

\begin{lstlisting}[language=bash]
openclaw status
openclaw models status
openclaw channels status
openclaw logs --follow
\end{lstlisting}


常见原因：

\begin{itemize}
  \item 模型认证未在 \textbf{Gateway 网关主机}上加载（检查 \texttt{models status}）。
  \item 渠道配对/允许列表阻止回复（检查渠道配置 + 日志）。
  \item WebChat/仪表板打开但没有正确的令牌。
\end{itemize}


如果你在远程，确认隧道/Tailscale 连接正常且 Gateway 网关 WebSocket 可达。

文档：\href{/channels}{渠道}、\href{/gateway/troubleshooting}{故障排除}、\href{/gateway/remote}{远程访问}。

\subsubsection{"Disconnected from gateway: no reason"——怎么办}


这通常意味着 UI 丢失了 WebSocket 连接。检查：

\begin{enumerate}
  \item Gateway 网关在运行吗？\texttt{openclaw gateway status}
  \item Gateway 网关健康吗？\texttt{openclaw status}
  \item UI 有正确的令牌吗？\texttt{openclaw dashboard}
  \item 如果是远程，隧道/Tailscale 链接正常吗？
\end{enumerate}


然后跟踪日志：

\begin{lstlisting}[language=bash]
openclaw logs --follow
\end{lstlisting}


文档：\href{/web/dashboard}{仪表板}、\href{/gateway/remote}{远程访问}、\href{/gateway/troubleshooting}{故障排除}。

\subsubsection{Telegram setMyCommands 因网络错误失败，应该检查什么}


从日志和渠道状态开始：

\begin{lstlisting}[language=bash]
openclaw channels status
openclaw channels logs --channel telegram
\end{lstlisting}


如果你在 VPS 上或代理后面，确认出站 HTTPS 被允许且 DNS 正常工作。
如果 Gateway 网关在远程，确保你在 Gateway 网关主机上查看日志。

文档：\href{/channels/telegram}{Telegram}、\href{/channels/troubleshooting}{渠道故障排除}。

\subsubsection{TUI 没有输出，应该检查什么}


首先确认 Gateway 网关可达且智能体可以运行：

\begin{lstlisting}[language=bash]
openclaw status
openclaw models status
openclaw logs --follow
\end{lstlisting}


在 TUI 中，使用 \texttt{/status} 查看当前状态。如果你期望在聊天渠道中收到回复，确保投递已启用（\texttt{/deliver on}）。

文档：\href{/web/tui}{TUI}、\href{/tools/slash-commands}{斜杠命令}。

\subsubsection{如何完全停止然后启动 Gateway 网关如果你安装了服务：}


\begin{lstlisting}[language=bash]
openclaw gateway stop
openclaw gateway start
\end{lstlisting}


这会停止/启动\textbf{受监管的服务}（macOS 上的 launchd，Linux 上的 systemd）。
当 Gateway 网关作为守护进程在后台运行时使用此命令。

如果你在前台运行，用 Ctrl‑C 停止，然后：

\begin{lstlisting}[language=bash]
openclaw gateway run
\end{lstlisting}


文档：\href{/gateway}{Gateway 网关服务运维手册}。

\subsubsection{通俗解释：openclaw gateway restart 与 openclaw gateway}


\begin{itemize}
  \item \texttt{openclaw gateway restart}：重启\textbf{后台服务}（launchd/systemd）。
  \item \texttt{openclaw gateway}：在这个终端会话中\textbf{前台}运行 Gateway 网关。
\end{itemize}


如果你安装了服务，使用 Gateway 网关命令。想要一次性前台运行时使用 \texttt{openclaw gateway}。

\subsubsection{出现故障时获取更多详情的最快方法是什么}


使用 \texttt{--verbose} 启动 Gateway 网关以获取更多控制台详情。然后检查日志文件中的渠道认证、模型路由和 RPC 错误。

\subsection{媒体与附件}


\subsubsection{我的 Skills 生成了图片/PDF，但什么都没发送}


智能体的出站附件必须包含 \texttt{MEDIA:} 行（独占一行）。参阅 \href{/start/openclaw}{OpenClaw 助手设置}和 \href{/tools/agent-send}{Agent send}。

CLI 发送：

\begin{lstlisting}[language=bash]
openclaw message send --target +15555550123 --message "Here you go" --media /path/to/file.png
\end{lstlisting}


还要检查：

\begin{itemize}
  \item 目标渠道支持出站媒体且未被允许列表阻止。
  \item 文件在提供商的大小限制内（图片会调整到最大 2048px）。
\end{itemize}


参阅\href{/nodes/images}{图片}。

\subsection{安全与访问控制}


\subsubsection{将 OpenClaw 暴露给入站私信安全吗}


将入站私信视为不可信输入。默认设计旨在降低风险：

\begin{itemize}
  \item 支持私信的渠道上的默认行为是\textbf{配对}：
  \item 未知发送者会收到配对码；机器人不处理他们的消息。
  \item 批准方式：\texttt{openclaw pairing approve  }
  \item 每个渠道的待处理请求上限为 \textbf{3 个}；如果没收到代码，检查 \texttt{openclaw pairing list }。
  \item 公开开放私信需要明确选择加入（\texttt{dmPolicy: "open"} 且允许列表 \texttt{"*"}）。
\end{itemize}


运行 \texttt{openclaw doctor} 以发现有风险的私信策略。

\subsubsection{提示注入只对公开机器人有影响吗}


不是。提示注入是关于\textbf{不可信内容}，不仅仅是谁能给机器人发私信。如果你的助手读取外部内容（网络搜索/抓取、浏览器页面、邮件、文档、附件、粘贴的日志），这些内容可能包含试图劫持模型的指令。即使\textbf{你是唯一的发送者}，这也可能发生。

最大的风险是在启用工具时：模型可能被诱导泄露上下文或代表你调用工具。通过以下方式减少影响范围：

\begin{itemize}
  \item 使用只读或禁用工具的“阅读器”智能体来总结不可信内容
  \item 对启用工具的智能体关闭 \texttt{web\_search} / \texttt{web\_fetch} / \texttt{browser}
  \item 沙箱隔离和严格的工具允许列表
\end{itemize}


详情：\href{/gateway/security}{安全}。

\subsubsection{我的机器人应该有自己的邮箱、GitHub 账户或电话号码吗}


是的，对于大多数设置来说。用独立的账户和电话号码隔离机器人可以在出问题时减少影响范围。这也使得轮换凭据或撤销访问更容易，而不影响你的个人账户。

从小处开始。只授予你实际需要的工具和账户的访问权限，以后需要时再扩展。

文档：\href{/gateway/security}{安全}、\href{/channels/pairing}{配对}。

\subsubsection{我能让它自主管理我的短信吗？这安全吗}


我们\textbf{不建议}完全自主管理你的个人消息。最安全的模式是：

\begin{itemize}
  \item 将私信保持在\textbf{配对模式}或严格的允许列表中。
  \item 如果你希望它代表你发消息，使用\textbf{独立的号码或账户}。
  \item 让它起草，然后\textbf{发送前批准}。
\end{itemize}


如果你想实验，在专用账户上进行并保持隔离。参阅\href{/gateway/security}{安全}。

\subsubsection{个人助理任务可以使用更便宜的模型吗}


可以，\textbf{如果}智能体仅用于聊天且输入是可信的。较小的模型更容易受到指令劫持，因此避免将它们用于启用工具的智能体或读取不可信内容时。如果你必须使用较小的模型，锁定工具并在沙箱中运行。参阅\href{/gateway/security}{安全}。

\subsubsection{我在 Telegram 中运行了 /start 但没收到配对码}


配对码\textbf{仅在}未知发送者向机器人发消息且 \texttt{dmPolicy: "pairing"} 启用时发送。\texttt{/start} 本身不会生成代码。

检查待处理请求：

\begin{lstlisting}[language=bash]
openclaw pairing list telegram
\end{lstlisting}


如果你想立即获得访问权限，将你的发送者 ID 加入允许列表或为该账户设置 \texttt{dmPolicy: "open"}。

\subsubsection{WhatsApp：会给我的联系人发消息吗？配对如何工作}


不会。WhatsApp 的默认私信策略是\textbf{配对}。未知发送者只会收到配对码，他们的消息\textbf{不会被处理}。OpenClaw 只回复它收到的聊天或你明确触发的发送。

批准配对：

\begin{lstlisting}[language=bash]
openclaw pairing approve whatsapp <code>
\end{lstlisting}


列出待处理请求：

\begin{lstlisting}[language=bash]
openclaw pairing list whatsapp
\end{lstlisting}


向导电话号码提示：它用于设置你的\textbf{允许列表/所有者}，以便你自己的私信被允许。它不用于自动发送。如果你在个人 WhatsApp 号码上运行，使用该号码并启用 \texttt{channels.whatsapp.selfChatMode}。

\subsection{聊天命令、中止任务和“停不下来”}


\subsubsection{如何阻止内部系统消息显示在聊天中}


大多数内部或工具消息只在该会话启用了 \textbf{verbose} 或 \textbf{reasoning} 时才出现。

在你看到它的聊天中修复：

\begin{lstlisting}
/verbose off
/reasoning off
\end{lstlisting}


如果仍然嘈杂，检查控制 UI 中的会话设置并将 verbose 设为\textbf{继承}。同时确认你没有使用在配置中将 \texttt{verboseDefault} 设为 \texttt{on} 的机器人配置文件。

文档：\href{/tools/thinking}{思考和详细输出}、\href{/gateway/security\#reasoning--verbose-output-in-groups}{安全}。

\subsubsection{如何停止/取消正在运行的任务}


发送以下任一内容作为\textbf{独立消息}（不带斜杠）：

\begin{lstlisting}
stop
abort
esc
wait
exit
interrupt
\end{lstlisting}


这些是中止触发器（不是斜杠命令）。

对于后台进程（来自 exec 工具），你可以要求智能体运行：

\begin{lstlisting}
process action:kill sessionId:XXX
\end{lstlisting}


斜杠命令概览：参阅\href{/tools/slash-commands}{斜杠命令}。

大多数命令必须作为以 \texttt{/} 开头的\textbf{独立}消息发送，但一些快捷方式（如 \texttt{/status}）对允许列表中的发送者也支持内联使用。

\subsubsection{如何从 Telegram 发送 Discord 消息？（"Cross-context messaging denied"）}


OpenClaw 默认阻止\textbf{跨提供商}消息。如果工具调用绑定到 Telegram，除非你明确允许，否则不会发送到 Discord。

为智能体启用跨提供商消息：

\begin{lstlisting}[language=json]
{
  agents: {
    defaults: {
      tools: {
        message: {
          crossContext: {
            allowAcrossProviders: true,
            marker: { enabled: true, prefix: "[from {channel}] " },
          },
        },
      },
    },
  },
}
\end{lstlisting}


编辑配置后重启 Gateway 网关。如果你只想为单个智能体设置，将其放在 \texttt{agents.list[].tools.message} 下。

\subsubsection{为什么感觉机器人“忽略”了快速连发的消息}


队列模式控制新消息如何与正在进行的运行交互。使用 \texttt{/queue} 更改模式：

\begin{itemize}
  \item \texttt{steer} - 新消息重定向当前任务
  \item \texttt{followup} - 逐条处理消息
  \item \texttt{collect} - 批量消息并回复一次（默认）
  \item \texttt{steer-backlog} - 立即转向，然后处理积压
  \item \texttt{interrupt} - 中止当前运行并重新开始
\end{itemize}


你可以为 followup 模式添加选项如 \texttt{debounce:2s cap:25 drop:summarize}。

\subsection{从截图/聊天记录中准确回答问题}


\textbf{问：“使用 API 密钥时 Anthropic 的默认模型是什么？”}

\textbf{答：} 在 OpenClaw 中，凭据和模型选择是分开的。设置 \texttt{ANTHROPIC\_API\_KEY}（或在认证配置文件中存储 Anthropic API 密钥）启用认证，但实际的默认模型是你在 \texttt{agents.defaults.model.primary} 中配置的（例如 \texttt{anthropic/claude-sonnet-4-5} 或 \texttt{anthropic/claude-opus-4-5}）。如果你看到 \texttt{No credentials found for profile "anthropic:default"}，意味着 Gateway 网关在正在运行的智能体的预期 \texttt{auth-profiles.json} 中找不到 Anthropic 凭据。

\hrulefill


仍然卡住？在 \href{https://discord.com/invite/clawd}{Discord} 中提问或发起 \href{https://github.com/openclaw/openclaw/discussions}{GitHub 讨论}。


\clearpage

% === 源文件: help/index.md ===
\section{帮助}


如果你想要快速“解决问题”的流程，从这里开始：

\begin{itemize}
  \item \textbf{故障排除：}\href{/help/troubleshooting}{从这里开始}
  \item \textbf{安装完整性检查（Node/npm/PATH）：}\href{/install\#nodejs--npm-path-sanity}{安装}
  \item \textbf{Gateway 网关问题：}\href{/gateway/troubleshooting}{Gateway 网关故障排除}
  \item \textbf{日志：}\href{/logging}{日志记录} 和 \href{/gateway/logging}{Gateway 网关日志记录}
  \item \textbf{修复：}\href{/gateway/doctor}{Doctor}
\end{itemize}


如果你在寻找概念性问题（不是“出问题了”）：

\begin{itemize}
  \item \href{/help/faq}{常见问题（概念）}
\end{itemize}



\clearpage

% === 源文件: help/scripts.md ===
\section{脚本}


\texttt{scripts/} 目录包含用于本地工作流和运维任务的辅助脚本。
当任务明确与某个脚本相关时使用这些脚本；否则优先使用 CLI。

\subsection{约定}


\begin{itemize}
  \item 除非在文档或发布检查清单中引用，否则脚本为\textbf{可选}。
  \item 当 CLI 接口存在时优先使用（例如：认证监控使用 \texttt{openclaw models status --check}）。
  \item 假定脚本与特定主机相关；在新机器上运行前请先阅读脚本内容。
\end{itemize}


\subsection{认证监控脚本}


认证监控脚本的文档请参阅：
\href{/automation/auth-monitoring}{/automation/auth-monitoring}

\subsection{添加脚本时}


\begin{itemize}
  \item 保持脚本专注且有文档说明。
  \item 在相关文档中添加简短条目（如果缺少则创建一个）。
\end{itemize}



\clearpage

% === 源文件: help/testing.md ===
\section{测试}


OpenClaw 包含三个 Vitest 测试套件（单元/集成、端到端、实时）以及一小组 Docker 运行器。

本文档是一份"我们如何测试"的指南：

\begin{itemize}
  \item 每个套件覆盖什么（以及它刻意\textit{不}覆盖什么）
  \item 常见工作流程应运行哪些命令（本地、推送前、调试）
  \item 实时测试如何发现凭证并选择模型/提供商
  \item 如何为现实中的模型/提供商问题添加回归测试
\end{itemize}


\subsection{快速开始}


日常使用：

\begin{itemize}
  \item 完整检查（推送前的预期流程）：\texttt{pnpm build \&\& pnpm check \&\& pnpm test}
\end{itemize}


当你修改测试或需要额外的信心时：

\begin{itemize}
  \item 覆盖率检查：\texttt{pnpm test:coverage}
  \item 端到端套件：\texttt{pnpm test:e2e}
\end{itemize}


调试真实提供商/模型时（需要真实凭证）：

\begin{itemize}
  \item 实时套件（模型 + Gateway 网关工具/图像探测）：\texttt{pnpm test:live}
\end{itemize}


提示：当你只需要一个失败用例时，建议使用下文描述的允许列表环境变量来缩小实时测试范围。

\subsection{测试套件（在哪里运行什么）}


可以将这些套件理解为"逐渐增强的真实性"（以及逐渐增加的不稳定性/成本）：

\subsubsection{单元/集成测试（默认）}


\begin{itemize}
  \item 命令：\texttt{pnpm test}
  \item 配置：\texttt{vitest.config.ts}
  \item 文件：\texttt{src/**/*.test.ts}
  \item 范围：
  \item 纯单元测试
  \item 进程内集成测试（Gateway 网关认证、路由、工具、解析、配置）
  \item 已知问题的确定性回归测试
  \item 预期：
  \item 在 CI 中运行
  \item 不需要真实密钥
  \item 应该快速且稳定
\end{itemize}


\subsubsection{端到端测试（Gateway 网关冒烟测试）}


\begin{itemize}
  \item 命令：\texttt{pnpm test:e2e}
  \item 配置：\texttt{vitest.e2e.config.ts}
  \item 文件：\texttt{src/**/*.e2e.test.ts}
  \item 范围：
  \item 多实例 Gateway 网关端到端行为
  \item WebSocket/HTTP 接口、节点配对和较重的网络操作
  \item 预期：
  \item 在 CI 中运行（当在流水线中启用时）
  \item 不需要真实密钥
  \item 比单元测试有更多活动部件（可能较慢）
\end{itemize}


\subsubsection{实时测试（真实提供商 + 真实模型）}


\begin{itemize}
  \item 命令：\texttt{pnpm test:live}
  \item 配置：\texttt{vitest.live.config.ts}
  \item 文件：\texttt{src/**/*.live.test.ts}
  \item 默认：通过 \texttt{pnpm test:live} \textbf{启用}（设置 \texttt{OPENCLAW\_LIVE\_TEST=1}）
  \item 范围：
  \item "这个提供商/模型用真实凭证\textit{今天}实际能工作吗？"
  \item 捕获提供商格式变更、工具调用怪癖、认证问题和速率限制行为
  \item 预期：
  \item 设计上不适合 CI 稳定运行（真实网络、真实提供商策略、配额、故障）
  \item 花费金钱/使用速率限制
  \item 建议运行缩小范围的子集而非"全部"
  \item 实时运行会加载 \texttt{\textasciitilde{}/.profile} 以获取缺失的 API 密钥
  \item Anthropic 密钥轮换：设置 \texttt{OPENCLAW\_LIVE\_ANTHROPIC\_KEYS="sk-...,sk-..."}（或 \texttt{OPENCLAW\_LIVE\_ANTHROPIC\_KEY=sk-...}）或多个 \texttt{ANTHROPIC\_API\_KEY*} 变量；测试会在遇到速率限制时重试
\end{itemize}


\subsection{我应该运行哪个套件？}


使用这个决策表：

\begin{itemize}
  \item 编辑逻辑/测试：运行 \texttt{pnpm test}（如果改动较大，加上 \texttt{pnpm test:coverage}）
  \item 涉及 Gateway 网关网络/WS 协议/配对：加上 \texttt{pnpm test:e2e}
  \item 调试"我的机器人挂了"/提供商特定故障/工具调用：运行缩小范围的 \texttt{pnpm test:live}
\end{itemize}


\subsection{实时测试：模型冒烟测试（配置文件密钥）}


实时测试分为两层，以便隔离故障：

\begin{itemize}
  \item "直接模型"告诉我们提供商/模型是否能用给定的密钥正常响应。
  \item "Gateway 网关冒烟测试"告诉我们完整的 Gateway 网关 + 智能体管道是否对该模型正常工作（会话、历史记录、工具、沙箱策略等）。
\end{itemize}


\subsubsection{第一层：直接模型补全（无 Gateway 网关）}


\begin{itemize}
  \item 测试：\texttt{src/agents/models.profiles.live.test.ts}
  \item 目标：
  \item 枚举发现的模型
  \item 使用 \texttt{getApiKeyForModel} 选择你有凭证的模型
  \item 每个模型运行一个小型补全（以及需要时的针对性回归测试）
  \item 如何启用：
  \item \texttt{pnpm test:live}（或直接调用 Vitest 时使用 \texttt{OPENCLAW\_LIVE\_TEST=1}）
  \item 设置 \texttt{OPENCLAW\_LIVE\_MODELS=modern}（或 \texttt{all}，modern 的别名）以实际运行此套件；否则会跳过以保持 \texttt{pnpm test:live} 专注于 Gateway 网关冒烟测试
  \item 如何选择模型：
  \item \texttt{OPENCLAW\_LIVE\_MODELS=modern} 运行现代允许列表（Opus/Sonnet/Haiku 4.5、GPT-5.x + Codex、Gemini 3、GLM 4.7、MiniMax M2.1、Grok 4）
  \item \texttt{OPENCLAW\_LIVE\_MODELS=all} 是现代允许列表的别名
  \item 或 \texttt{OPENCLAW\_LIVE\_MODELS="openai/gpt-5.2,anthropic/claude-opus-4-5,..."}（逗号分隔的允许列表）
  \item 如何选择提供商：
  \item \texttt{OPENCLAW\_LIVE\_PROVIDERS="google,google-antigravity,google-gemini-cli"}（逗号分隔的允许列表）
  \item 密钥来源：
  \item 默认：配置文件存储和环境变量回退
  \item 设置 \texttt{OPENCLAW\_LIVE\_REQUIRE\_PROFILE\_KEYS=1} 以强制\textbf{仅使用配置文件存储}
  \item 为什么存在这个测试：
  \item 将"提供商 API 损坏/密钥无效"与"Gateway 网关智能体管道损坏"分离
  \item 包含小型、隔离的回归测试（例如：OpenAI Responses/Codex Responses 推理重放 + 工具调用流程）
\end{itemize}


\subsubsection{第二层：Gateway 网关 + 开发智能体冒烟测试（"@openclaw"实际做的事）}


\begin{itemize}
  \item 测试：\texttt{src/gateway/gateway-models.profiles.live.test.ts}
  \item 目标：
  \item 启动一个进程内 Gateway 网关
  \item 创建/修补一个 \texttt{agent:dev:*} 会话（每次运行覆盖模型）
  \item 遍历有密钥的模型并断言：
  \item "有意义的"响应（无工具）
  \item 真实的工具调用工作正常（读取探测）
  \item 可选的额外工具探测（执行+读取探测）
  \item OpenAI 回归路径（仅工具调用 → 后续）保持工作
  \item 探测详情（以便你能快速解释故障）：
  \item \texttt{read} 探测：测试在工作区写入一个随机数文件，要求智能体 \texttt{read} 它并回显随机数。
  \item \texttt{exec+read} 探测：测试要求智能体 \texttt{exec} 将随机数写入临时文件，然后 \texttt{read} 回来。
  \item 图像探测：测试附加一个生成的 PNG（猫 + 随机代码），期望模型返回 \texttt{cat }。
  \item 实现参考：\texttt{src/gateway/gateway-models.profiles.live.test.ts} 和 \texttt{src/gateway/live-image-probe.ts}。
  \item 如何启用：
  \item \texttt{pnpm test:live}（或直接调用 Vitest 时使用 \texttt{OPENCLAW\_LIVE\_TEST=1}）
  \item 如何选择模型：
  \item 默认：现代允许列表（Opus/Sonnet/Haiku 4.5、GPT-5.x + Codex、Gemini 3、GLM 4.7、MiniMax M2.1、Grok 4）
  \item \texttt{OPENCLAW\_LIVE\_GATEWAY\_MODELS=all} 是现代允许列表的别名
  \item 或设置 \texttt{OPENCLAW\_LIVE\_GATEWAY\_MODELS="provider/model"}（或逗号分隔列表）来缩小范围
  \item 如何选择提供商（避免"OpenRouter 全部"）：
  \item \texttt{OPENCLAW\_LIVE\_GATEWAY\_PROVIDERS="google,google-antigravity,google-gemini-cli,openai,anthropic,zai,minimax"}（逗号分隔的允许列表）
  \item 工具 + 图像探测在此实时测试中始终开启：
  \item \texttt{read} 探测 + \texttt{exec+read} 探测（工具压力测试）
  \item 当模型声明支持图像输入时运行图像探测
  \item 流程（高层次）：
  \item 测试生成一个带有"CAT"+ 随机代码的小型 PNG（\texttt{src/gateway/live-image-probe.ts}）
  \item 通过 \texttt{agent} \texttt{attachments: [\{ mimeType: "image/png", content: "" \}]} 发送
  \item Gateway 网关将附件解析为 \texttt{images[]}（\texttt{src/gateway/server-methods/agent.ts} + \texttt{src/gateway/chat-attachments.ts}）
  \item 嵌入式智能体将多模态用户消息转发给模型
  \item 断言：回复包含 \texttt{cat} + 代码（OCR 容差：允许轻微错误）
\end{itemize}


提示：要查看你的机器上可以测试什么（以及确切的 \texttt{provider/model} ID），运行：

\begin{lstlisting}[language=bash]
openclaw models list
openclaw models list --json
\end{lstlisting}


\subsection{实时测试：Anthropic 设置令牌冒烟测试}


\begin{itemize}
  \item 测试：\texttt{src/agents/anthropic.setup-token.live.test.ts}
  \item 目标：验证 Claude Code CLI 设置令牌（或粘贴的设置令牌配置文件）能完成 Anthropic 提示。
  \item 启用：
  \item \texttt{pnpm test:live}（或直接调用 Vitest 时使用 \texttt{OPENCLAW\_LIVE\_TEST=1}）
  \item \texttt{OPENCLAW\_LIVE\_SETUP\_TOKEN=1}
  \item 令牌来源（选择一个）：
  \item 配置文件：\texttt{OPENCLAW\_LIVE\_SETUP\_TOKEN\_PROFILE=anthropic:setup-token-test}
  \item 原始令牌：\texttt{OPENCLAW\_LIVE\_SETUP\_TOKEN\_VALUE=sk-ant-oat01-...}
  \item 模型覆盖（可选）：
  \item \texttt{OPENCLAW\_LIVE\_SETUP\_TOKEN\_MODEL=anthropic/claude-opus-4-5}
\end{itemize}


设置示例：

\begin{lstlisting}[language=bash]
openclaw models auth paste-token --provider anthropic --profile-id anthropic:setup-token-test
OPENCLAW_LIVE_SETUP_TOKEN=1 OPENCLAW_LIVE_SETUP_TOKEN_PROFILE=anthropic:setup-token-test pnpm test:live src/agents/anthropic.setup-token.live.test.ts
\end{lstlisting}


\subsection{实时测试：CLI 后端冒烟测试（Claude Code CLI 或其他本地 CLI）}


\begin{itemize}
  \item 测试：\texttt{src/gateway/gateway-cli-backend.live.test.ts}
  \item 目标：使用本地 CLI 后端验证 Gateway 网关 + 智能体管道，而不影响你的默认配置。
  \item 启用：
  \item \texttt{pnpm test:live}（或直接调用 Vitest 时使用 \texttt{OPENCLAW\_LIVE\_TEST=1}）
  \item \texttt{OPENCLAW\_LIVE\_CLI\_BACKEND=1}
  \item 默认值：
  \item 模型：\texttt{claude-cli/claude-sonnet-4-5}
  \item 命令：\texttt{claude}
  \item 参数：\texttt{["-p","--output-format","json","--dangerously-skip-permissions"]}
  \item 覆盖（可选）：
  \item \texttt{OPENCLAW\_LIVE\_CLI\_BACKEND\_MODEL="claude-cli/claude-opus-4-5"}
  \item \texttt{OPENCLAW\_LIVE\_CLI\_BACKEND\_MODEL="codex-cli/gpt-5.2-codex"}
  \item \texttt{OPENCLAW\_LIVE\_CLI\_BACKEND\_COMMAND="/full/path/to/claude"}
  \item \texttt{OPENCLAW\_LIVE\_CLI\_BACKEND\_ARGS='["-p","--output-format","json","--permission-mode","bypassPermissions"]'}
  \item \texttt{OPENCLAW\_LIVE\_CLI\_BACKEND\_CLEAR\_ENV='["ANTHROPIC\_API\_KEY","ANTHROPIC\_API\_KEY\_OLD"]'}
  \item \texttt{OPENCLAW\_LIVE\_CLI\_BACKEND\_IMAGE\_PROBE=1} 发送真实图像附件（路径注入到提示中）。
  \item \texttt{OPENCLAW\_LIVE\_CLI\_BACKEND\_IMAGE\_ARG="--image"} 将图像文件路径作为 CLI 参数传递而非提示注入。
  \item \texttt{OPENCLAW\_LIVE\_CLI\_BACKEND\_IMAGE\_MODE="repeat"}（或 \texttt{"list"}）控制设置 \texttt{IMAGE\_ARG} 时如何传递图像参数。
  \item \texttt{OPENCLAW\_LIVE\_CLI\_BACKEND\_RESUME\_PROBE=1} 发送第二轮并验证恢复流程。
  \item \texttt{OPENCLAW\_LIVE\_CLI\_BACKEND\_DISABLE\_MCP\_CONFIG=0} 保持 Claude Code CLI MCP 配置启用（默认使用临时空文件禁用 MCP 配置）。
\end{itemize}


示例：

\begin{lstlisting}[language=bash]
OPENCLAW_LIVE_CLI_BACKEND=1 \
  OPENCLAW_LIVE_CLI_BACKEND_MODEL="claude-cli/claude-sonnet-4-5" \
  pnpm test:live src/gateway/gateway-cli-backend.live.test.ts
\end{lstlisting}


\subsubsection{推荐的实时测试配方}


缩小范围的显式允许列表最快且最不易出错：

\begin{itemize}
  \item 单个模型，直接测试（无 Gateway 网关）：
  \item \texttt{OPENCLAW\_LIVE\_MODELS="openai/gpt-5.2" pnpm test:live src/agents/models.profiles.live.test.ts}
\end{itemize}


\begin{itemize}
  \item 单个模型，Gateway 网关冒烟测试：
  \item \texttt{OPENCLAW\_LIVE\_GATEWAY\_MODELS="openai/gpt-5.2" pnpm test:live src/gateway/gateway-models.profiles.live.test.ts}
\end{itemize}


\begin{itemize}
  \item 跨多个提供商的工具调用：
  \item \texttt{OPENCLAW\_LIVE\_GATEWAY\_MODELS="openai/gpt-5.2,anthropic/claude-opus-4-5,google/gemini-3-flash-preview,zai/glm-4.7,minimax/minimax-m2.1" pnpm test:live src/gateway/gateway-models.profiles.live.test.ts}
\end{itemize}


\begin{itemize}
  \item Google 专项（Gemini API 密钥 + Antigravity）：
  \item Gemini（API 密钥）：\texttt{OPENCLAW\_LIVE\_GATEWAY\_MODELS="google/gemini-3-flash-preview" pnpm test:live src/gateway/gateway-models.profiles.live.test.ts}
  \item Antigravity（OAuth）：\texttt{OPENCLAW\_LIVE\_GATEWAY\_MODELS="google-antigravity/claude-opus-4-6-thinking,google-antigravity/gemini-3-pro-high" pnpm test:live src/gateway/gateway-models.profiles.live.test.ts}
\end{itemize}


注意：

\begin{itemize}
  \item \texttt{google/...} 使用 Gemini API（API 密钥）。
  \item \texttt{google-antigravity/...} 使用 Antigravity OAuth 桥接（Cloud Code Assist 风格的智能体端点）。
  \item \texttt{google-gemini-cli/...} 使用你机器上的本地 Gemini CLI（独立的认证 + 工具怪癖）。
  \item Gemini API 与 Gemini CLI：
  \item API：OpenClaw 通过 HTTP 调用 Google 托管的 Gemini API（API 密钥/配置文件认证）；这是大多数用户说的"Gemini"。
  \item CLI：OpenClaw 调用本地 \texttt{gemini} 二进制文件；它有自己的认证，行为可能不同（流式传输/工具支持/版本差异）。
\end{itemize}


\subsection{实时测试：模型矩阵（我们覆盖什么）}


没有固定的"CI 模型列表"（实时测试是可选的），但这些是建议在有密钥的开发机器上定期覆盖的\textbf{推荐}模型。

\subsubsection{现代冒烟测试集（工具调用 + 图像）}


这是我们期望保持工作的"常用模型"运行：

\begin{itemize}
  \item OpenAI（非 Codex）：\texttt{openai/gpt-5.2}（可选：\texttt{openai/gpt-5.1}）
  \item OpenAI Codex：\texttt{openai-codex/gpt-5.2}（可选：\texttt{openai-codex/gpt-5.2-codex}）
  \item Anthropic：\texttt{anthropic/claude-opus-4-5}（或 \texttt{anthropic/claude-sonnet-4-5}）
  \item Google（Gemini API）：\texttt{google/gemini-3-pro-preview} 和 \texttt{google/gemini-3-flash-preview}（避免较旧的 Gemini 2.x 模型）
  \item Google（Antigravity）：\texttt{google-antigravity/claude-opus-4-6-thinking} 和 \texttt{google-antigravity/gemini-3-flash}
  \item Z.AI（GLM）：\texttt{zai/glm-4.7}
  \item MiniMax：\texttt{minimax/minimax-m2.1}
\end{itemize}


运行带工具 + 图像的 Gateway 网关冒烟测试：
\texttt{OPENCLAW\_LIVE\_GATEWAY\_MODELS="openai/gpt-5.2,openai-codex/gpt-5.2,anthropic/claude-opus-4-5,google/gemini-3-pro-preview,google/gemini-3-flash-preview,google-antigravity/claude-opus-4-6-thinking,google-antigravity/gemini-3-flash,zai/glm-4.7,minimax/minimax-m2.1" pnpm test:live src/gateway/gateway-models.profiles.live.test.ts}

\subsubsection{基线：工具调用（Read + 可选 Exec）}


每个提供商系列至少选择一个：

\begin{itemize}
  \item OpenAI：\texttt{openai/gpt-5.2}（或 \texttt{openai/gpt-5-mini}）
  \item Anthropic：\texttt{anthropic/claude-opus-4-5}（或 \texttt{anthropic/claude-sonnet-4-5}）
  \item Google：\texttt{google/gemini-3-flash-preview}（或 \texttt{google/gemini-3-pro-preview}）
  \item Z.AI（GLM）：\texttt{zai/glm-4.7}
  \item MiniMax：\texttt{minimax/minimax-m2.1}
\end{itemize}


可选的额外覆盖（锦上添花）：

\begin{itemize}
  \item xAI：\texttt{xai/grok-4}（或最新可用版本）
  \item Mistral：\texttt{mistral/}…（选择一个你已启用的"工具"能力模型）
  \item Cerebras：\texttt{cerebras/}…（如果你有访问权限）
  \item LM Studio：\texttt{lmstudio/}…（本地；工具调用取决于 API 模式）
\end{itemize}


\subsubsection{视觉：图像发送（附件 → 多模态消息）}


在 \texttt{OPENCLAW\_LIVE\_GATEWAY\_MODELS} 中至少包含一个支持图像的模型（Claude/Gemini/OpenAI 视觉能力变体等）以测试图像探测。

\subsubsection{聚合器/替代 Gateway 网关}


如果你启用了密钥，我们也支持通过以下方式测试：

\begin{itemize}
  \item OpenRouter：\texttt{openrouter/...}（数百个模型；使用 \texttt{openclaw models scan} 查找支持工具+图像的候选模型）
  \item OpenCode Zen：\texttt{opencode/...}（通过 \texttt{OPENCODE\_API\_KEY} / \texttt{OPENCODE\_ZEN\_API\_KEY} 认证）
\end{itemize}


如果你有凭证/配置，可以在实时矩阵中包含更多提供商：

\begin{itemize}
  \item 内置：\texttt{openai}、\texttt{openai-codex}、\texttt{anthropic}、\texttt{google}、\texttt{google-vertex}、\texttt{google-antigravity}、\texttt{google-gemini-cli}、\texttt{zai}、\texttt{openrouter}、\texttt{opencode}、\texttt{xai}、\texttt{groq}、\texttt{cerebras}、\texttt{mistral}、\texttt{github-copilot}
  \item 通过 \texttt{models.providers}（自定义端点）：\texttt{minimax}（云/API），以及任何 OpenAI/Anthropic 兼容代理（LM Studio、vLLM、LiteLLM 等）
\end{itemize}


提示：不要试图在文档中硬编码"所有模型"。权威列表是你机器上 \texttt{discoverModels(...)} 返回的内容 + 可用的密钥。

\subsection{凭证（绝不提交）}


实时测试以与 CLI 相同的方式发现凭证。实际含义：

\begin{itemize}
  \item 如果 CLI 能工作，实时测试应该能找到相同的密钥。
  \item 如果实时测试说"无凭证"，用调试 \texttt{openclaw models list}/模型选择相同的方式调试。
\end{itemize}


\begin{itemize}
  \item 配置文件存储：\texttt{\textasciitilde{}/.openclaw/credentials/}（首选；测试中"配置文件密钥"的含义）
  \item 配置：\texttt{\textasciitilde{}/.openclaw/openclaw.json}（或 \texttt{OPENCLAW\_CONFIG\_PATH}）
\end{itemize}


如果你想依赖环境变量密钥（例如在 \texttt{\textasciitilde{}/.profile} 中导出的），在 \texttt{source \textasciitilde{}/.profile} 后运行本地测试，或使用下面的 Docker 运行器（它们可以将 \texttt{\textasciitilde{}/.profile} 挂载到容器中）。

\subsection{Deepgram 实时测试（音频转录）}


\begin{itemize}
  \item 测试：\texttt{src/media-understanding/providers/deepgram/audio.live.test.ts}
  \item 启用：\texttt{DEEPGRAM\_API\_KEY=... DEEPGRAM\_LIVE\_TEST=1 pnpm test:live src/media-understanding/providers/deepgram/audio.live.test.ts}
\end{itemize}


\subsection{Docker 运行器（可选的"在 Linux 中工作"检查）}


这些在仓库 Docker 镜像内运行 \texttt{pnpm test:live}，挂载你的本地配置目录和工作区（如果挂载了 \texttt{\textasciitilde{}/.profile} 则会加载它）：

\begin{itemize}
  \item 直接模型：\texttt{pnpm test:docker:live-models}（脚本：\texttt{scripts/test-live-models-docker.sh}）
  \item Gateway 网关 + 开发智能体：\texttt{pnpm test:docker:live-gateway}（脚本：\texttt{scripts/test-live-gateway-models-docker.sh}）
  \item 新手引导向导（TTY，完整脚手架）：\texttt{pnpm test:docker:onboard}（脚本：\texttt{scripts/e2e/onboard-docker.sh}）
  \item Gateway 网关网络（两个容器，WS 认证 + 健康检查）：\texttt{pnpm test:docker:gateway-network}（脚本：\texttt{scripts/e2e/gateway-network-docker.sh}）
  \item 插件（自定义扩展加载 + 注册表冒烟测试）：\texttt{pnpm test:docker:plugins}（脚本：\texttt{scripts/e2e/plugins-docker.sh}）
\end{itemize}


有用的环境变量：

\begin{itemize}
  \item \texttt{OPENCLAW\_CONFIG\_DIR=...}（默认：\texttt{\textasciitilde{}/.openclaw}）挂载到 \texttt{/home/node/.openclaw}
  \item \texttt{OPENCLAW\_WORKSPACE\_DIR=...}（默认：\texttt{\textasciitilde{}/.openclaw/workspace}）挂载到 \texttt{/home/node/.openclaw/workspace}
  \item \texttt{OPENCLAW\_PROFILE\_FILE=...}（默认：\texttt{\textasciitilde{}/.profile}）挂载到 \texttt{/home/node/.profile} 并在运行测试前加载
  \item \texttt{OPENCLAW\_LIVE\_GATEWAY\_MODELS=...} / \texttt{OPENCLAW\_LIVE\_MODELS=...} 用于缩小运行范围
  \item \texttt{OPENCLAW\_LIVE\_REQUIRE\_PROFILE\_KEYS=1} 确保凭证来自配置文件存储（而非环境变量）
\end{itemize}


\subsection{文档完整性检查}


文档编辑后运行文档检查：\texttt{pnpm docs:list}。

\subsection{离线回归测试（CI 安全）}


这些是没有真实提供商的"真实管道"回归测试：

\begin{itemize}
  \item Gateway 网关工具调用（模拟 OpenAI，真实 Gateway 网关 + 智能体循环）：\texttt{src/gateway/gateway.tool-calling.mock-openai.test.ts}
  \item Gateway 网关向导（WS \texttt{wizard.start}/\texttt{wizard.next}，写入配置 + 强制认证）：\texttt{src/gateway/gateway.wizard.e2e.test.ts}
\end{itemize}


\subsection{智能体可靠性评估（Skills）}


我们已经有一些 CI 安全的测试，它们的行为类似于"智能体可靠性评估"：

\begin{itemize}
  \item 通过真实 Gateway 网关 + 智能体循环的模拟工具调用（\texttt{src/gateway/gateway.tool-calling.mock-openai.test.ts}）。
  \item 验证会话连接和配置效果的端到端向导流程（\texttt{src/gateway/gateway.wizard.e2e.test.ts}）。
\end{itemize}


对于 Skills 仍然缺少的内容（参见 \href{/tools/skills}{Skills}）：

\begin{itemize}
  \item \textbf{决策：} 当 Skills 在提示中列出时，智能体是否选择正确的 skill（或避免不相关的）？
  \item \textbf{合规性：} 智能体是否在使用前读取 \texttt{SKILL.md} 并遵循所需的步骤/参数？
  \item \textbf{工作流契约：} 断言工具顺序、会话历史延续和沙箱边界的多轮场景。
\end{itemize}


未来的评估应该首先保持确定性：

\begin{itemize}
  \item 使用模拟提供商来断言工具调用 + 顺序、skill 文件读取和会话连接的场景运行器。
  \item 一小套专注于 skill 的场景（使用 vs 避免、门控、提示注入）。
  \item 可选的实时评估（可选的，环境变量门控），仅在 CI 安全套件就位后。
\end{itemize}


\subsection{添加回归测试（指导）}


当你修复在实时测试中发现的提供商/模型问题时：

\begin{itemize}
  \item 如果可能，添加 CI 安全的回归测试（模拟/存根提供商，或捕获确切的请求形状转换）
  \item 如果它本质上是仅限实时的（速率限制、认证策略），保持实时测试范围小且通过环境变量可选
  \item 优先针对能捕获问题的最小层：
  \item 提供商请求转换/重放问题 → 直接模型测试
  \item Gateway 网关会话/历史/工具管道问题 → Gateway 网关实时冒烟测试或 CI 安全的 Gateway 网关模拟测试
\end{itemize}



\clearpage

% === 源文件: help/troubleshooting.md ===
\section{故障排除}


\subsection{最初的六十秒}


按顺序运行这些命令：

\begin{lstlisting}[language=bash]
openclaw status
openclaw status --all
openclaw gateway probe
openclaw logs --follow
openclaw doctor
\end{lstlisting}


如果 Gateway 网关可达，进行深度探测：

\begin{lstlisting}[language=bash]
openclaw status --deep
\end{lstlisting}


\subsection{常见的“它坏了”情况}


\subsubsection{openclaw: command not found}


几乎总是 Node/npm PATH 问题。从这里开始：

\begin{itemize}
  \item \href{/install\#nodejs--npm-path-sanity}{安装（Node/npm PATH 安装完整性检查）}
\end{itemize}


\subsubsection{安装程序失败（或你需要完整日志）}


以详细模式重新运行安装程序以查看完整跟踪和 npm 输出：

\begin{lstlisting}[language=bash]
curl -fsSL https://openclaw.ai/install.sh | bash -s -- --verbose
\end{lstlisting}


对于 beta 安装：

\begin{lstlisting}[language=bash]
curl -fsSL https://openclaw.ai/install.sh | bash -s -- --beta --verbose
\end{lstlisting}


你也可以设置 \texttt{OPENCLAW\_VERBOSE=1} 代替标志。

\subsubsection{Gateway 网关“unauthorized”、无法连接或持续重连}


\begin{itemize}
  \item \href{/gateway/troubleshooting}{Gateway 网关故障排除}
  \item \href{/gateway/authentication}{Gateway 网关认证}
\end{itemize}


\subsubsection{控制 UI 在 HTTP 上失败（需要设备身份）}


\begin{itemize}
  \item \href{/gateway/troubleshooting}{Gateway 网关故障排除}
  \item \href{/web/control-ui\#insecure-http}{控制 UI}
\end{itemize}


\subsubsection{docs.openclaw.ai 显示 SSL 错误（Comcast/Xfinity）}


一些 Comcast/Xfinity 连接通过 Xfinity Advanced Security 阻止 \texttt{docs.openclaw.ai}。
禁用 Advanced Security 或将 \texttt{docs.openclaw.ai} 添加到允许列表，然后重试。

\begin{itemize}
  \item Xfinity Advanced Security 帮助：https://www.xfinity.com/support/articles/using-xfinity-xfi-advanced-security
  \item 快速检查：尝试移动热点或 VPN 以确认这是 ISP 级别的过滤
\end{itemize}


\subsubsection{服务显示运行中，但 RPC 探测失败}


\begin{itemize}
  \item \href{/gateway/troubleshooting}{Gateway 网关故障排除}
  \item \href{/gateway/background-process}{后台进程/服务}
\end{itemize}


\subsubsection{模型/认证失败（速率限制、账单、“all models failed”）}


\begin{itemize}
  \item \href{/cli/models}{模型}
  \item \href{/concepts/oauth}{OAuth / 认证概念}
\end{itemize}


\subsubsection{/model 显示 model not allowed}


这通常意味着 \texttt{agents.defaults.models} 配置为允许列表。当它非空时，只能选择那些提供商/模型键。

\begin{itemize}
  \item 检查允许列表：\texttt{openclaw config get agents.defaults.models}
  \item 添加你想要的模型（或清除允许列表）然后重试 \texttt{/model}
  \item 使用 \texttt{/models} 浏览允许的提供商/模型
\end{itemize}


\subsubsection{提交问题时}


粘贴一份安全报告：

\begin{lstlisting}[language=bash]
openclaw status --all
\end{lstlisting}


如果可以的话，包含来自 \texttt{openclaw logs --follow} 的相关日志尾部。


\clearpage
